\documentclass[11pt, a4paper]{article}

% --- UNIVERSAL PREAMBLE BLOCK ---
\usepackage[a4paper, top=2.5cm, bottom=2.5cm, left=2cm, right=2cm]{geometry}
\usepackage{fontspec}

\usepackage[english, bidi=basic, provide=*]{babel}

\babelprovide[import, onchar=ids fonts]{english}

% Set default/Latin font to Sans Serif in the main (rm) slot
\babelfont{rm}{Noto Sans}

% --- ADDITIONAL PACKAGES ---
\usepackage{amssymb}    % For checkboxes \square
\usepackage{booktabs}   % For professional tables
\usepackage{longtable}  % For multi-page tracking table
\usepackage{enumitem}   % For customized lists
\usepackage{xcolor}     % For colored highlights
\usepackage[colorlinks=true, linkcolor=blue]{hyperref} % For internal links
\usepackage{tikz}
\usetikzlibrary{arrows.meta, positioning, shapes.geometric, backgrounds, fit}
\usepackage{listings}
\lstdefinestyle{code}{
    basicstyle=\ttfamily\footnotesize,
    breaklines=true,
    frame=single,
    backgroundcolor=\color{gray!8},
    rulecolor=\color{gray!40},
    columns=fullflexible,
    keepspaces=true,
    showstringspaces=false
}
\lstset{style=code}

% --- CUSTOM COMMANDS ---
\newcommand{\project}[1]{\vspace{1mm}\noindent\textbf{\textcolor{blue}{Hands-on Project:}} #1}
\newcommand{\usecase}[1]{\vspace{1mm}\noindent\textbf{\textcolor{teal}{Use Case \& Architecture:}} #1}
\newcommand{\capstone}[2]{
    \vspace{3mm}
    \noindent\fbox{
        \parbox{0.97\textwidth}{
            \textbf{\textcolor{red}{CAPSTONE PROJECT - #1:}} #2
        }
    }
    \vspace{3mm}
}
\newcommand{\dayitem}[2]{\item \textbf{#1:} #2}
% Rubric box: capstone acceptance criteria
\newcommand{\rubric}[1]{
    \vspace{1mm}
    \noindent\fbox{
        \parbox{0.97\textwidth}{
            \textbf{\textcolor{orange}{Acceptance Criteria:}} #1
        }
    }
    \vspace{2mm}
}
% Weekly curated resources
\newcommand{\resources}[1]{\vspace{1mm}\noindent\textbf{\textcolor{violet}{Resources:}} #1}
% Weekly self-check questions
\newcommand{\selfcheck}[1]{\vspace{1mm}\noindent\textbf{\textcolor{purple}{Self-check:}} #1}
% Production hardening note
\newcommand{\prodnote}[1]{
    \vspace{1mm}
    \noindent\fbox{
        \parbox{0.97\textwidth}{
            \textbf{\textcolor{brown}{Production Note:}} \textit{#1}
        }
    }
    \vspace{1mm}
}

\title{\Huge \textbf{Databricks Master Data Engineer}\\ \Large 6-Month Exhaustive Training \& Certification Program}
\author{Advanced Architecture, Hands-on Labs, and MLOps}
\date{}

\begin{document}

\maketitle

\begin{abstract}
\noindent This 24-week program is designed to transform you from a practitioner to a Master Data Engineer on the Databricks Data Intelligence Platform. It progressively builds from lakehouse and Delta Lake foundations to advanced Unity Catalog governance, declarative pipelines, structured streaming, Mosaic AI, and production-grade CI/CD, integrating daily theory, weekly service-specific mini-projects, architectural use case analyses, and massive monthly Capstone projects. A tracking timetable is included at the end.
\end{abstract}

\tableofcontents
\newpage

\section*{Prerequisites \& How to Use This Syllabus}
\begin{itemize}[label=-]
    \item \textbf{Prerequisite skills:} intermediate Python (PySpark helpful), solid SQL, Git workflows (branching, pull requests), Linux/CLI basics, and REST API fundamentals.
    \item \textbf{Weekly cadence:} approximately 10--12 hours per week --- Monday to Wednesday for theory, Thursday for the hands-on project, Friday for the architecture use case.
    \item \textbf{Deliverable expectation:} every capstone must produce a git repository, an architecture diagram, a cost estimate (in DBUs and cloud infrastructure), and a security review checklist, per the Acceptance Criteria boxes that follow each capstone.
    \item \textbf{Assessment:} the 24-week tracking timetable at the end of this document doubles as the assessment checklist --- check off each week only when its project and deliverables are complete.
\end{itemize}

\noindent\textbf{Environment setup checklist:}
\begin{itemize}[label=-]
    \item Install Python 3.x, Git, Docker Desktop, and VS Code (with the Databricks extension).
    \item Install the Databricks CLI v0.200+ (\texttt{pip install databricks-cli} or the native installer) and \texttt{databricks-sdk-py}; verify with \texttt{databricks auth profiles} and \texttt{databricks workspace list /}.
    \item Create a free Databricks trial workspace (or use Databricks Community Edition / an existing organizational workspace) with guardrails: a dedicated catalog/schema for labs, cluster auto-termination set to 30 minutes or less, and a budget alert on the cloud account.
    \item Clone the course repository structure: \texttt{data/}, \texttt{src/}, \texttt{tests/}, \texttt{bundles/}, \texttt{docs/}.
    \item Configure credentials via a CLI profile with OAuth (\texttt{databricks auth login}) or a personal access token stored in an environment variable --- never hardcode tokens in code or notebooks.
\end{itemize}

% ==========================================
% MILESTONE 1
% ==========================================
\section{Milestone 1: Lakehouse Foundations (Month 1)}
\textit{Objective: Master the Databricks workspace, Delta Lake storage, the medallion architecture, and Apache Spark fundamentals.}
\par\noindent\textbf{Learning outcomes:} navigate and administer a Databricks workspace, explain Delta Lake's transaction log and ACID guarantees, design Bronze/Silver/Gold data flows, and write idiomatic PySpark against managed Delta tables.

\subsection*{Week 1: Databricks Platform \& Workspace Fundamentals}
\par\noindent\textbf{Outcomes:} explain the Databricks control plane vs.\ data plane split, navigate workspaces, notebooks, and repos, choose between all-purpose and jobs compute, and authenticate the Databricks CLI and Python SDK.
\begin{itemize}[label=-]
    \dayitem{Monday}{Platform architecture: control plane, data plane, serverless vs.\ classic compute, workspace objects.}
    \dayitem{Tuesday}{Notebooks, Repos (Git folders), magic commands, and \texttt{dbutils} utilities.}
    \dayitem{Wednesday}{Compute: all-purpose clusters, job clusters, policies, pools, and auto-termination; the Databricks CLI and \texttt{databricks-sdk-py}.}
    \dayitem{Thursday}{\project{Provision a trial workspace, configure the Databricks CLI with an OAuth profile, create an all-purpose cluster with 30-minute auto-termination, clone a Git repo into a Git folder, and run a smoke-test notebook via \texttt{databricks jobs submit} using the SDK.}}
    \dayitem{Friday}{\usecase{Design the workspace and compute layout for a 200-person data organization spanning three business units with dev/staging/prod isolation.}}
\end{itemize}
\resources{Databricks documentation (\url{https://docs.databricks.com/}); Data Engineering Zoomcamp (\url{https://github.com/DataTalksClub/data-engineering-zoomcamp}).}
\selfcheck{\begin{itemize}[label=-, leftmargin=1.5em, itemsep=0pt]
    \item Which resources live in the control plane and which in the data plane?
    \item Why are job clusters cheaper than all-purpose clusters for scheduled pipelines?
    \item What does a cluster policy restrict, and who benefits from that restriction?
    \item How does a Git folder (Repo) differ from a workspace notebook folder?
    \item Why prefer OAuth (\texttt{databricks auth login}) over a long-lived personal access token?
\end{itemize}}

\prodnote{Set auto-termination on every non-serverless cluster before your first real workload: an idle all-purpose cluster left running over a weekend is the single most common first-month bill shock. Enforce it with a cluster policy that caps \texttt{autotermination\_minutes} and deny users the \texttt{unpinned} permission on shared clusters.}

\subsection*{Week 2: Delta Lake Foundations}
\par\noindent\textbf{Outcomes:} explain the Delta transaction log, ACID semantics, and schema enforcement/evolution, use time travel and \texttt{DESCRIBE HISTORY}, and recover from accidental writes with \texttt{RESTORE}.
\begin{itemize}[label=-]
    \dayitem{Monday}{Delta Lake architecture: Parquet data files, the \texttt{\_delta\_log} transaction log, and atomic commits.}
    \dayitem{Tuesday}{ACID guarantees, optimistic concurrency, isolation levels, and \texttt{MERGE INTO} semantics.}
    \dayitem{Wednesday}{Schema enforcement and evolution, constraints, generated columns, time travel, and \texttt{RESTORE}.}
    \dayitem{Thursday}{\project{Create Delta tables from CSV seed data, force a schema mismatch and observe enforcement, evolve the schema deliberately with \texttt{mergeSchema}, delete a partition ``by accident,'' then recover it with time travel and \texttt{RESTORE}; document every step from \texttt{DESCRIBE HISTORY}.}}
    \dayitem{Friday}{\usecase{Design the landing and recovery strategy for a financial-services lake that must prove to auditors exactly what the data looked like at any point in the last 90 days.}}
\end{itemize}
\resources{Delta Lake documentation (\url{https://docs.delta.io/}); Delta Lake VLDB paper (\url{https://www.vldb.org/pvldb/vol13/p3411-armbrust.pdf}).}
\selfcheck{\begin{itemize}[label=-, leftmargin=1.5em, itemsep=0pt]
    \item What makes a Delta commit atomic even though it spans many Parquet files?
    \item Which isolation level does Delta use for concurrent writes to one table?
    \item \texttt{mergeSchema} vs.\ \texttt{overwriteSchema}: when is each appropriate?
    \item What does \texttt{VACUUM} permanently destroy, and why does the 7-day default exist?
    \item Why is \texttt{MERGE INTO} preferred over delete-then-insert for upserts?
\end{itemize}}

\prodnote{Time travel is a recovery tool with an expiry date: \texttt{VACUUM} removes the data files that older versions reference. Never lower retention below your longest recovery objective, and treat any \texttt{RETAIN} under 168 hours as a change-control event requiring sign-off.}

\subsection*{Week 3: The Medallion Architecture (Bronze/Silver/Gold)}
\par\noindent\textbf{Outcomes:} justify the Bronze/Silver/Gold layering, design data-quality gates between layers, model a Gold star schema, and name catalogs/schemas/tables for a multi-team lakehouse.
\begin{itemize}[label=-]
    \dayitem{Monday}{Medallion rationale: progressive refinement, idempotent reprocessing, and the cost of skipping layers.}
    \dayitem{Tuesday}{Bronze patterns: append-only raw capture, metadata columns (\texttt{\_ingested\_at}, \texttt{\_source\_file}), and schema-on-read trade-offs.}
    \dayitem{Wednesday}{Silver cleansing and conformance, Gold dimensional modeling, and Unity Catalog naming (\texttt{catalog.schema.table}).}
    \dayitem{Thursday}{\project{Build a three-layer medallion flow for the Olist e-commerce dataset: land raw CSVs into Bronze with ingestion metadata, cleanse and deduplicate into Silver with quarantine for rejected rows, and publish a Gold \texttt{fact\_orders} plus two dimensions; reconcile row counts across layers.}}
    \dayitem{Friday}{\usecase{Architect the medallion lakehouse for a retail chain that must reprocess two years of history after a business-rule change without double-counting revenue.}}
\end{itemize}
\resources{Databricks medallion architecture guide (\url{https://www.databricks.com/glossary/medallion-architecture}); Data Engineer Handbook (\url{https://github.com/DataExpert-io/data-engineer-handbook}).}
\selfcheck{\begin{itemize}[label=-, leftmargin=1.5em, itemsep=0pt]
    \item Why keep Bronze append-only even when storage is ``cheap''?
    \item What belongs in a quarantine table instead of being silently dropped?
    \item How does the medallion layout make a re-run idempotent?
    \item Why model Gold as a star schema rather than a wide denormalized table?
    \item Which metadata columns prove lineage from Gold back to the raw file?
\end{itemize}}

\subsection*{Week 4: Apache Spark on Databricks}
\par\noindent\textbf{Outcomes:} explain Spark's driver/executor model and lazy evaluation, read the query plan from the Spark UI, use DataFrame transformations vs.\ actions, and partition data deliberately with \texttt{repartition} and \texttt{coalesce}.
\begin{itemize}[label=-]
    \dayitem{Monday}{Spark architecture: driver, executors, DAG, stages, tasks, and lazy evaluation.}
    \dayitem{Tuesday}{DataFrame API: transformations vs.\ actions, Catalyst optimizer, and the query plan (\texttt{explain}).}
    \dayitem{Wednesday}{Partitions, shuffles, skew, caching, and the Spark UI on Databricks.}
    \dayitem{Thursday}{\project{Write a PySpark notebook that ingests 10M NYC Taxi rows, runs aggregations with deliberate narrow and wide transformations, then use the Spark UI to find the shuffle stage, identify skewed partitions, and fix them with salting or \texttt{repartition}.}}
    \dayitem{Friday}{\usecase{Diagnose a nightly job that grew from 40 minutes to 6 hours as data grew 10x: produce a tuning plan grounded in the Spark UI evidence you would collect.}}
\end{itemize}
\resources{Apache Spark documentation (\url{https://spark.apache.org/docs/latest/}); Databricks runtime release notes (\url{https://docs.databricks.com/release-notes/runtime/}).}
\selfcheck{\begin{itemize}[label=-, leftmargin=1.5em, itemsep=0pt]
    \item Why does lazy evaluation let Catalyst optimize your job?
    \item What triggers a shuffle, and why is it expensive?
    \item \texttt{repartition} vs.\ \texttt{coalesce}: which can increase partition count?
    \item What does a long tail in the Spark UI stage summary usually indicate?
    \item When does \texttt{cache()} hurt more than it helps?
\end{itemize}}

\capstone{Milestone 1}{Build a Retail Medallion Lakehouse. Ingest the Olist dataset into a Bronze schema with ingestion metadata, cleanse and deduplicate into Silver with a quarantine table and row-count reconciliation, publish a Gold star schema (\texttt{fact\_orders}, \texttt{dim\_customer}, \texttt{dim\_product}, \texttt{dim\_date}), and demonstrate time-travel recovery of a deliberately corrupted Silver table.}
\rubric{\begin{itemize}[label=-, leftmargin=1.5em, itemsep=0pt]
    \item End-to-end pipeline reruns from a clean clone following the repo README.
    \item Architecture diagram committed (draw.io/TikZ/PNG) showing all three medallion layers and data-quality gates.
    \item At least 3 automated data-quality checks (row-count reconciliation, uniqueness of \texttt{order\_id}, quarantine accounting) with a failing-case demonstration.
    \item \texttt{DESCRIBE HISTORY} evidence of the corruption and \texttt{RESTORE} recovery.
    \item Monthly cost estimate (DBU/hour for the cluster class used, plus cloud storage) with 2 optimization decisions documented (e.g., auto-termination, job clusters).
    \item Security review: workspace-level access notes, no hardcoded tokens, PII columns in \texttt{dim\_customer} identified.
\end{itemize}}

\subsection*{Milestone 1 Capstone: Reference Solution Sketch}
\begin{itemize}[label=-, itemsep=0pt]
    \item \textbf{Architecture:} volume-landed CSVs $\rightarrow$ Bronze (append-only Delta with \texttt{\_ingested\_at}) $\rightarrow$ Silver (typed, deduplicated, quarantined rejects) $\rightarrow$ Gold star schema --- decouples raw evidence from curated analytics.
    \item \textbf{Steps:} create catalog \texttt{training} with schemas \texttt{bronze}, \texttt{silver}, \texttt{gold}; upload Olist CSVs to a Unity Catalog volume; write parameterized PySpark notebooks per layer; add reconciliation notebook asserting \texttt{accepted + rejected = input}; corrupt Silver with a bad \texttt{DELETE}; recover with \texttt{RESTORE TO VERSION AS OF}.
    \item \textbf{Hardest part --- idempotent Silver merge:}
\begin{lstlisting}
MERGE INTO silver.orders t
USING bronze_orders_deduped s
ON t.order_id = s.order_id
WHEN MATCHED THEN UPDATE SET *
WHEN NOT MATCHED THEN INSERT *
\end{lstlisting}
    \item \textbf{Pitfalls:} running \texttt{VACUUM} before the recovery demo destroys the versions you need; a non-parameterized notebook forces code edits per re-run; forgetting to stop the all-purpose cluster leaks DBUs overnight.
\end{itemize}

\begin{figure}[htbp]
\centering
\begin{tikzpicture}[
    node distance=0.55cm and 1.0cm,
    box/.style={draw, rounded corners, fill=blue!10, align=center, font=\small, minimum height=0.9cm, inner sep=4pt},
    side/.style={draw, rounded corners, fill=brown!10, align=center, font=\footnotesize, inner sep=3pt},
    arr/.style={-{Stealth}, thick}
]
\node[box] (vol) {UC Volume\\(raw CSVs)};
\node[box, right=of vol, fill=brown!12] (bronze) {Bronze\\append-only Delta};
\node[box, right=of bronze, fill=gray!15] (silver) {Silver\\typed + deduped};
\node[box, right=of silver, fill=yellow!18] (gold) {Gold\\star schema};
\node[side, above=0.5cm of silver] (quar) {Quarantine table};
\node[side, below=0.5cm of silver] (hist) {DESCRIBE HISTORY\\+ RESTORE};
\node[box, right=of gold] (bi) {SQL Warehouse\\consumers};
\draw[arr] (vol) -- (bronze);
\draw[arr] (bronze) -- (silver);
\draw[arr] (silver) -- (gold);
\draw[arr] (gold) -- (bi);
\draw[arr, dashed] (silver) -- (quar);
\draw[arr, dashed] (hist) -- (silver);
\end{tikzpicture}
\caption{Milestone 1 reference architecture: medallion lakehouse with quality gates and time-travel recovery}
\end{figure}

% ==========================================
% MILESTONE 2
% ==========================================
\section{Milestone 2: Warehousing \& Analytics (Month 2)}
\textit{Objective: Optimize Delta storage, serve SQL analytics with warehouses, and govern data with Unity Catalog.}
\par\noindent\textbf{Learning outcomes:} compact and cluster Delta tables for selective queries, operate SQL warehouses (serverless and classic), build AI/BI dashboards and Genie spaces, and implement catalog-level governance with Unity Catalog.

\subsection*{Week 5: Delta Lake Optimization (OPTIMIZE, Z-ORDER, Liquid Clustering)}
\par\noindent\textbf{Outcomes:} diagnose the small-file problem, run \texttt{OPTIMIZE} with \texttt{ZORDER BY}, explain liquid clustering and when it supersedes partitioning/Z-ORDER, and schedule \texttt{VACUUM} safely.
\begin{itemize}[label=-]
    \dayitem{Monday}{File layout and data skipping: file-level statistics, the small-file problem, and bloom filters.}
    \dayitem{Tuesday}{\texttt{OPTIMIZE} compaction, \texttt{ZORDER BY} multi-dimensional clustering, and auto-optimize vs.\ scheduled jobs.}
    \dayitem{Wednesday}{Liquid clustering (\texttt{CLUSTER BY}), predictive optimization, deletion vectors, and \texttt{VACUUM} retention.}
    \dayitem{Thursday}{\project{Create a 50M-row events table with thousands of small files, benchmark a selective query, then compact with \texttt{OPTIMIZE ZORDER BY}, convert a second copy to liquid clustering, benchmark again, and record file counts and timings in a results table.}}
    \dayitem{Friday}{\usecase{Remediate a lakehouse whose nightly streaming sink produced 2 million tiny files: design the compaction, clustering, and write-cadence plan, and defend liquid clustering vs.\ partition + Z-ORDER.}}
\end{itemize}
\resources{Delta Lake optimization documentation (\url{https://docs.databricks.com/delta/optimize.html}); liquid clustering announcement and docs (\url{https://docs.databricks.com/delta/clustering.html}).}
\selfcheck{\begin{itemize}[label=-, leftmargin=1.5em, itemsep=0pt]
    \item Why do many small files slow queries even when total bytes are unchanged?
    \item Which columns make good \texttt{ZORDER BY} candidates?
    \item What does liquid clustering let you change without rewriting the table?
    \item What do deletion vectors avoid rewriting on \texttt{DELETE}?
    \item Why should \texttt{OPTIMIZE} not run against the table a stream is writing at the same moment?
\end{itemize}}

\prodnote{Prefer liquid clustering for new large tables: it avoids the partition-column decision you will regret in year two, allows \texttt{CLUSTER BY} changes without a full rewrite, and composes with predictive optimization so maintenance runs itself on serverless. Keep \texttt{ZORDER} for legacy tables you cannot yet migrate.}

\subsection*{Week 6: Databricks SQL \& SQL Warehouses}
\par\noindent\textbf{Outcomes:} size and operate SQL warehouses (serverless, pro, classic), write warehouse-grade SQL against Unity Catalog tables, use materialized views and result caching, and control concurrency and cost.
\begin{itemize}[label=-]
    \dayitem{Monday}{Databricks SQL architecture: SQL warehouses, serverless vs.\ pro vs.\ classic, Photon, and the query profile.}
    \dayitem{Tuesday}{Analytical SQL: window functions, Delta time travel in SQL, \texttt{QUALIFY}, and query federation basics.}
    \dayitem{Wednesday}{Materialized views and streaming tables in DBSQL, result caching, warehouse scaling, and monitoring with system tables.}
    \dayitem{Thursday}{\project{Load the Gold star schema from Month 1 behind a serverless SQL warehouse, build a materialized view for daily revenue, capture the query profile before/after, and configure min/max clusters plus auto-stop; document cost per query from \texttt{system.billing.usage}.}}
    \dayitem{Friday}{\usecase{Design the warehouse topology for a company where 500 analysts run morning dashboards while data engineers run backfills on the same tables without contention.}}
\end{itemize}
\resources{Databricks SQL documentation (\url{https://docs.databricks.com/sql/}); Photon engine overview (\url{https://www.databricks.com/product/photon}).}
\selfcheck{\begin{itemize}[label=-, leftmargin=1.5em, itemsep=0pt]
    \item What does Photon accelerate and why does it matter for SQL workloads?
    \item When does a materialized view beat a scheduled Gold table?
    \item Serverless vs.\ classic SQL warehouse: what do you give up and gain?
    \item How does max-clusters scaling absorb the 9 a.m.\ dashboard burst?
    \item Which system table shows who spent what on SQL compute?
\end{itemize}}

\subsection*{Week 7: AI/BI Dashboards \& Genie}
\par\noindent\textbf{Outcomes:} build AI/BI dashboards on governed Gold data, configure Genie spaces with trusted instructions, curate datasets and measures for self-service, and publish dashboards on a schedule.
\begin{itemize}[label=-]
    \dayitem{Monday}{AI/BI dashboards: datasets, visualizations, parameters, and cross-filters.}
    \dayitem{Tuesday}{AI/BI Genie: natural-language analytics, instructions, example queries, and trusted assets.}
    \dayitem{Wednesday}{Semantic curation: certified Gold tables, measures, lineage into dashboards, and scheduled refresh/sharing.}
    \dayitem{Thursday}{\project{Build an executive AI/BI dashboard on \texttt{gold.fact\_orders} with three parameterized pages, then create a Genie space over the same schema with instructions that correctly answer ``revenue by state last quarter''; capture the generated SQL and fix one wrong answer with a certified measure.}}
    \dayitem{Friday}{\usecase{Design the self-service analytics layer for a retailer where business users ask Genie ad-hoc questions while finance requires certified, reconciled numbers only.}}
\end{itemize}
\resources{AI/BI dashboards documentation (\url{https://docs.databricks.com/dashboards/}); Genie documentation (\url{https://docs.databricks.com/genie/}).}
\selfcheck{\begin{itemize}[label=-, leftmargin=1.5em, itemsep=0pt]
    \item Why must Genie spaces sit on Gold, never Bronze?
    \item What do Genie ``instructions'' change about generated SQL?
    \item Dashboard dataset vs.\ direct table reference: what is the trade-off?
    \item How do you certify a number so finance can trust the dashboard?
    \item What does a scheduled dashboard subscription need that ad-hoc viewing does not?
\end{itemize}}

\subsection*{Week 8: Unity Catalog Foundations \& Data Governance}
\par\noindent\textbf{Outcomes:} explain the metastore$\rightarrow$catalog$\rightarrow$schema$\rightarrow$table hierarchy, grant least-privilege access with \texttt{GRANT}, use volumes for file data, and audit access with system tables.
\begin{itemize}[label=-]
    \dayitem{Monday}{Unity Catalog architecture: metastores, catalogs, schemas, tables, volumes, and the three-level namespace.}
    \dayitem{Tuesday}{Privilege model: \texttt{GRANT}/\texttt{REVOKE}, ownership, groups, and row/column-level security with dynamic views.}
    \dayitem{Wednesday}{Volumes vs.\ external locations, managed vs.\ external tables, audit logging, and data discovery.}
    \dayitem{Thursday}{\project{Create a catalog with \texttt{bronze}/\texttt{silver}/\texttt{gold} schemas, register an external volume, grant a read-only analyst group Gold-only access, build a dynamic view masking customer PII columns, and prove the mask by querying as the restricted principal.}}
    \dayitem{Friday}{\usecase{Design governance for a healthcare lakehouse where analysts query Gold, engineers write Silver, only pipelines write Bronze, and every access must be auditable for HIPAA.}}
\end{itemize}
\resources{Unity Catalog documentation (\url{https://docs.databricks.com/data-governance/unity-catalog/}); Data Governance with Unity Catalog whitepaper (\url{https://www.databricks.com/resources}).}
\selfcheck{\begin{itemize}[label=-, leftmargin=1.5em, itemsep=0pt]
    \item What does a metastore bind together, and how many per region?
    \item Managed vs.\ external table: who owns the lifecycle of the files?
    \item Why use a dynamic view for masking instead of a copied sanitized table?
    \item Which privilege lets a group read Gold but nothing beneath it?
    \item Where does the audit record of ``who selected what'' live?
\end{itemize}}

\capstone{Milestone 2}{Build a Governed Analytics Lakehouse. Extend the Month 1 medallion lakehouse with optimized storage (liquid clustering on fact tables), a serverless SQL warehouse serving an AI/BI dashboard plus a Genie space, and Unity Catalog governance: analyst/engineer groups, dynamic-view PII masking, and audit evidence from system tables.}
\rubric{\begin{itemize}[label=-, leftmargin=1.5em, itemsep=0pt]
    \item Benchmark table showing before/after timings for \texttt{OPTIMIZE ZORDER BY} and liquid clustering on the same query.
    \item AI/BI dashboard with at least 3 pages and parameters; Genie space answering 5 business questions with documented instructions.
    \item Access matrix implemented (not just described): engineers read/write Silver, analysts read Gold via a dynamic masked view, proven with a test principal.
    \item Audit evidence: \texttt{system.access.audit} query showing a masked-column read by the analyst group.
    \item Cost estimate in DBUs for warehouse + pipeline compute with 2 optimizations (auto-stop, serverless) documented.
    \item Security review: no external-table sprawl, no hardcoded credentials, PII inventory table committed.
\end{itemize}}

\subsection*{Milestone 2 Capstone: Reference Solution Sketch}
\begin{itemize}[label=-, itemsep=0pt]
    \item \textbf{Architecture:} Gold Delta tables (liquid clustered) $\rightarrow$ serverless SQL warehouse $\rightarrow$ AI/BI dashboard + Genie space; Unity Catalog enforces group-based grants and a masking dynamic view; system tables provide the audit trail.
    \item \textbf{Steps:} \texttt{ALTER TABLE ... CLUSTER BY} on the fact table and re-benchmark; create warehouse with auto-stop 10 min; build dashboard datasets from certified views; create analyst/engineer groups and \texttt{GRANT}s; write \texttt{v\_customer\_masked} with \texttt{CASE WHEN is\_member('analysts')}; validate with a test user.
    \item \textbf{Hardest part --- the masking view:}
\begin{lstlisting}
CREATE OR REPLACE VIEW gold.v_customer_masked AS
SELECT customer_id,
  CASE WHEN is_member('pii_readers') THEN email
       ELSE sha2(email, 256) END AS email,
  region, created_at
FROM gold.dim_customer;
\end{lstlisting}
    \item \textbf{Pitfalls:} granting at catalog level when Gold-only was required; Genie instructions contradicting the certified measure; benchmarking before clustering and drawing the wrong conclusion.
\end{itemize}

\begin{figure}[htbp]
\centering
\begin{tikzpicture}[
    node distance=0.55cm and 0.9cm,
    box/.style={draw, rounded corners, fill=blue!10, align=center, font=\small, minimum height=0.9cm, inner sep=4pt},
    sec/.style={draw, rounded corners, fill=red!8, align=center, font=\footnotesize, inner sep=3pt},
    arr/.style={-{Stealth}, thick}
]
\node[box] (gold) {Gold Delta\\(liquid clustered)};
\node[box, right=of gold] (wh) {Serverless\\SQL Warehouse};
\node[box, right=of wh, yshift=0.9cm] (dash) {AI/BI\\Dashboard};
\node[box, right=of wh, yshift=-0.9cm] (genie) {Genie\\Space};
\node[sec, below=0.5cm of gold] (uc) {Unity Catalog\\GRANTs + masking view};
\node[sec, below=0.5cm of wh] (audit) {system.access.audit};
\draw[arr] (gold) -- (wh);
\draw[arr] (wh) -- (dash);
\draw[arr] (wh) -- (genie);
\draw[arr, dashed] (uc) -- (gold);
\draw[arr, dashed] (audit) -- (wh);
\end{tikzpicture}
\caption{Milestone 2 reference architecture: governed, optimized analytics serving layer}
\end{figure}

% ==========================================
% MILESTONE 3
% ==========================================
\section{Milestone 3: Batch ETL \& Orchestration (Month 3)}
\textit{Objective: Engineer production batch pipelines with PySpark, Delta Live Tables, Auto Loader, and Databricks Workflows.}
\par\noindent\textbf{Learning outcomes:} tune PySpark jobs, declare pipelines with DLT expectations, ingest incrementally with Auto Loader and CDC, and orchestrate multi-task jobs with Workflows.

\subsection*{Week 9: PySpark Deep Dive \& Performance Tuning}
\par\noindent\textbf{Outcomes:} choose between RDD, DataFrame, and SQL APIs deliberately, tune shuffle partitions and joins (broadcast vs.\ sort-merge), handle skew with salting, and profile with the Spark UI and Ganglia.
\begin{itemize}[label=-]
    \dayitem{Monday}{Catalyst and Tungsten under the hood: whole-stage codegen, memory management, and when to drop to RDDs (rarely).}
    \dayitem{Tuesday}{Joins deep dive: broadcast hash vs.\ sort-merge, AQE (adaptive query execution), and skew hints.}
    \dayitem{Wednesday}{Shuffle tuning: \texttt{spark.sql.shuffle.partitions}, salting, caching strategy, and UDF performance (pandas UDFs).}
    \dayitem{Thursday}{\project{Take the Week 4 job and tune it end-to-end: fix shuffle partition count, convert one join to broadcast, salt a skewed key, replace a Python UDF with a pandas UDF, and produce a before/after evidence table from Spark UI metrics.}}
    \dayitem{Friday}{\usecase{A partner's daily 2 TB join job misses its SLA by 3 hours: write the performance investigation report you would deliver, with the exact metrics and config changes.}}
\end{itemize}
\resources{Spark SQL performance tuning guide (\url{https://spark.apache.org/docs/latest/sql-performance-tuning.html}); The Definitive Guide to Apache Spark (Chambers \& Zaharia).}
\selfcheck{\begin{itemize}[label=-, leftmargin=1.5em, itemsep=0pt]
    \item What does AQE change at runtime that static planning cannot?
    \item When does a broadcast join become dangerous?
    \item Why is a Python UDF slower than a pandas UDF?
    \item How does salting fix skew, and what does it cost the other side of the join?
    \item Which Spark UI tab tells you a stage is shuffle-bound?
\end{itemize}}

\subsection*{Week 10: Delta Live Tables (Declarative Pipelines)}
\par\noindent\textbf{Outcomes:} define streaming tables and materialized views declaratively, enforce quality with expectations (\texttt{expect}, \texttt{expect\_or\_drop}, \texttt{expect\_or\_fail}), run triggered vs.\ continuous pipelines, and read the DLT event log.
\begin{itemize}[label=-]
    \dayitem{Monday}{DLT concepts: datasets, dependencies, declarative DAGs, and development vs.\ production modes.}
    \dayitem{Tuesday}{Expectations and data quality: warn/drop/fail semantics and quality metrics.}
    \dayitem{Wednesday}{Pipeline operations: triggered vs.\ continuous, full refresh, schema evolution handling, and the event log.}
    \dayitem{Thursday}{\project{Rebuild the Month 1 Bronze$\rightarrow$Silver$\rightarrow$Gold flow as a DLT pipeline with expectations dropping bad \texttt{order\_id}s and failing on null \texttt{customer\_id}s; inject a poison record, observe drop metrics, and query the event log for quality statistics.}}
    \dayitem{Friday}{\usecase{Choose DLT vs.\ hand-written Spark + Workflows for three real pipeline candidates at a logistics company, and defend the boundary you draw.}}
\end{itemize}
\resources{Delta Live Tables documentation (\url{https://docs.databricks.com/delta-live-tables/}); DLT expectations guide (\url{https://docs.databricks.com/delta-live-tables/expectations.html}).}
\selfcheck{\begin{itemize}[label=-, leftmargin=1.5em, itemsep=0pt]
    \item What does the DLT DAG give you that a notebook chain does not?
    \item \texttt{expect\_or\_drop} vs.\ \texttt{expect\_or\_fail}: pick one for a PII column and justify.
    \item Why is the event log itself a Delta table worth querying?
    \item When is a full refresh required even in an incremental pipeline?
    \item Continuous vs.\ triggered: which suits a nightly cost-sensitive load?
\end{itemize}}

\subsection*{Week 11: Auto Loader, Incremental Ingestion \& CDC}
\par\noindent\textbf{Outcomes:} ingest files incrementally with Auto Loader (\texttt{cloudFiles}) in directory and notification modes, rescue malformed records with \texttt{rescuedDataColumn}, apply CDC with \texttt{MERGE} and \texttt{apply\_changes} patterns, and manage schema inference/evolution.
\begin{itemize}[label=-]
    \dayitem{Monday}{Auto Loader architecture: cloudFiles source, directory listing vs.\ file notification mode, checkpoints, and exactly-once semantics.}
    \dayitem{Tuesday}{Schema inference, schema hints, evolution, and the rescued data column.}
    \dayitem{Wednesday}{CDC patterns: \texttt{MERGE INTO} upserts, SCD Type 1/2 with \texttt{AUTO CDC ... INTO} (formerly \texttt{apply\_changes}), and delete handling.}
    \dayitem{Thursday}{\project{Drop daily CSV batches into a UC volume, ingest them with an Auto Loader streaming query in triggered mode, simulate a schema change mid-stream plus a corrupted row, then apply a CDC feed to a Silver table with SCD Type 2 history for customers.}}
    \dayitem{Friday}{\usecase{Design ingestion for a payments platform receiving 500 small files/minute from partners plus a nightly CDC extract from the core banking database.}}
\end{itemize}
\resources{Auto Loader documentation (\url{https://docs.databricks.com/ingestion/auto-loader/}); SCD with DLT documentation (\url{https://docs.databricks.com/delta-live-tables/cdc.html}).}
\selfcheck{\begin{itemize}[label=-, leftmargin=1.5em, itemsep=0pt]
    \item Notification vs.\ directory-listing mode: when does the choice matter at scale?
    \item What makes Auto Loader exactly-once across restarts?
    \item What lands in \texttt{\_rescued\_data} and why is that better than failing?
    \item SCD Type 1 vs.\ Type 2: what does each answer about history?
    \item How does \texttt{SEQUENCE BY} resolve out-of-order CDC events?
\end{itemize}}

\prodnote{Auto Loader's checkpoint is the pipeline's memory: never delete it to ``fix'' a stuck stream --- you will reprocess or skip files. If you must reset, version the target table first or plan a full backfill with reconciliation.}

\subsection*{Week 12: Workflows, Jobs \& Pipeline Orchestration}
\par\noindent\textbf{Outcomes:} build multi-task jobs with dependencies, parameterize tasks, configure retries and alerts, orchestrate DLT pipelines and notebooks together, and use \texttt{dbutils.notebook.run} vs.\ task values correctly.
\begin{itemize}[label=-]
    \dayitem{Monday}{Jobs and tasks: DAGs of tasks, job clusters vs.\ shared clusters, and scheduling with cron.}
    \dayitem{Tuesday}{Parameters, task values, conditionals (\texttt{if/else} tasks), and for-each tasks.}
    \dayitem{Wednesday}{Reliability: retries, timeouts, alerts, repair runs, and monitoring with system tables.}
    \dayitem{Thursday}{\project{Orchestrate the full medallion flow as one job: Auto Loader ingestion task $\rightarrow$ DLT pipeline task $\rightarrow$ data-quality notebook $\rightarrow$ conditional alert task on failure, parameterized by \texttt{processing\_date}; force a failure and execute a repair run of only the failed tasks.}}
    \dayitem{Friday}{\usecase{Design the orchestration layer for 60 interdependent nightly pipelines across three teams with a shared 6 a.m.\ SLA, including backfill and failure-isolation strategy.}}
\end{itemize}
\resources{Databricks Workflows documentation (\url{https://docs.databricks.com/workflows/}); Jobs API reference (\url{https://docs.databricks.com/api/workspace/jobs}).}
\selfcheck{\begin{itemize}[label=-, leftmargin=1.5em, itemsep=0pt]
    \item Why give each task its own job cluster in a critical pipeline?
    \item Task values vs.\ job parameters: which is computed at runtime?
    \item What does a repair run re-execute, and why is that cheaper than a full rerun?
    \item Where should a data-quality gate sit relative to Gold publishing?
    \item How do \texttt{if/else} tasks replace external alerting glue?
\end{itemize}}

\capstone{Milestone 3}{Build an Orchestrated Batch ELT Platform. Ingest the full Olist dataset daily via Auto Loader into Bronze, transform with a DLT pipeline enforcing expectations into Silver/Gold, orchestrate with a multi-task Workflow (ingest $\rightarrow$ pipeline $\rightarrow$ quality gate $\rightarrow$ publish) parameterized by date, with retries, failure alerts, and a demonstrated repair run.}
\rubric{\begin{itemize}[label=-, leftmargin=1.5em, itemsep=0pt]
    \item One Workflow job runs the entire flow; a forced mid-pipeline failure is repaired (not fully rerun) with evidence.
    \item DLT expectations produce measurable quality metrics; one poison record demonstrably dropped and counted.
    \item Auto Loader handles a schema-evolving batch and a corrupted row via \texttt{\_rescued\_data} without manual repair.
    \item Pipeline is idempotent: a same-day rerun changes zero Gold row counts.
    \item Cost estimate per nightly run in DBUs, with job-cluster vs.\ shared-cluster decision documented.
    \item Security review: pipeline runs as a service principal, no tokens in notebooks, checkpoint locations governed.
\end{itemize}}

\subsection*{Milestone 3 Capstone: Reference Solution Sketch}
\begin{itemize}[label=-, itemsep=0pt]
    \item \textbf{Architecture:} partner files $\rightarrow$ Auto Loader (volume + checkpoint) $\rightarrow$ Bronze streaming table $\rightarrow$ DLT Silver with expectations $\rightarrow$ Gold publish task $\rightarrow$ conditional failure alert.
    \item \textbf{Steps:} create the volume and checkpoint schema; author the DLT pipeline; build the Workflow with task values passing \texttt{processing\_date}; add the quality-gate notebook asserting reconciliation; wire the \texttt{if/else} alert task; run, break, repair.
    \item \textbf{Hardest part --- the quality gate as a task:}
\begin{lstlisting}
bad = spark.sql("""
  SELECT COUNT(*) c FROM silver.orders
  WHERE order_id IS NULL OR total_amount < 0
""").first().c
if bad > 0:
    raise ValueError(f"Quality gate failed: {bad} bad rows")
dbutils.jobs.taskValues.set("quality", "passed")
\end{lstlisting}
    \item \textbf{Pitfalls:} deleting the Auto Loader checkpoint and double-ingesting; expectations logging but not wired to \texttt{expect\_or\_fail} where it matters; a shared cluster for all tasks turning one OOM into six failed tasks.
\end{itemize}

\begin{figure}[htbp]
\centering
\begin{tikzpicture}[
    node distance=0.55cm and 0.9cm,
    box/.style={draw, rounded corners, fill=blue!10, align=center, font=\small, minimum height=0.9cm, inner sep=4pt},
    gate/.style={draw, rounded corners, fill=orange!15, align=center, font=\footnotesize, inner sep=3pt},
    arr/.style={-{Stealth}, thick}
]
\node[box] (files) {Partner files\\(UC volume)};
\node[box, right=of files] (al) {Auto Loader\\(+ checkpoint)};
\node[box, right=of al] (dlt) {DLT pipeline\\expectations};
\node[gate, right=of dlt] (gate) {Quality gate\\notebook};
\node[box, right=of gate] (gold) {Gold publish};
\node[gate, below=0.5cm of gate] (alert) {if/else alert task};
\draw[arr] (files) -- (al);
\draw[arr] (al) -- (dlt);
\draw[arr] (dlt) -- (gate);
\draw[arr] (gate) -- (gold);
\draw[arr, dashed] (gate) -- (alert);
\end{tikzpicture}
\caption{Milestone 3 reference architecture: orchestrated, quality-gated batch ELT}
\end{figure}

% ==========================================
% MILESTONE 4
% ==========================================
\section{Milestone 4: Streaming \& Real-Time (Month 4)}
\textit{Objective: Build production streaming pipelines with Structured Streaming, stateful processing, and Kafka integration.}
\par\noindent\textbf{Learning outcomes:} reason in event time with windows and watermarks, guarantee exactly-once with checkpoints and idempotent sinks, integrate Kafka/Event Hubs/Kinesis, and architect real-time lakehouse systems.

\subsection*{Week 13: Structured Streaming Foundations}
\par\noindent\textbf{Outcomes:} explain the micro-batch execution model, read from streaming sources, write to Delta with checkpoints and triggers, and monitor a stream with \texttt{StreamingQueryListener} and the UI.
\begin{itemize}[label=-]
    \dayitem{Monday}{Structured Streaming model: the infinite table, micro-batches vs.\ continuous, and the Catalyst integration.}
    \dayitem{Tuesday}{Sources, sinks, output modes (append/update/complete), and triggers including \texttt{availableNow}.}
    \dayitem{Wednesday}{Checkpoints, offsets, recovery semantics, and stream monitoring.}
    \dayitem{Thursday}{\project{Build a streaming pipeline from a simulated clickstream (rate source or files dropped into a volume) into a Delta sink; run with three different triggers, kill and restart the query mid-stream, and prove from the checkpoint that no events were lost or duplicated.}}
    \dayitem{Friday}{\usecase{Design the streaming ingestion layer for an ad-tech platform receiving 50k events/second that must land queryable data within one minute.}}
\end{itemize}
\resources{Structured Streaming programming guide (\url{https://spark.apache.org/docs/latest/structured-streaming-programming-guide.html}); Databricks streaming documentation (\url{https://docs.databricks.com/structured-streaming/}).}
\selfcheck{\begin{itemize}[label=-, leftmargin=1.5em, itemsep=0pt]
    \item Why is a checkpoint directory mandatory for a production stream?
    \item Append vs.\ update vs.\ complete output mode: which supports aggregation to Delta?
    \item What does \texttt{availableNow} trigger give a streaming job?
    \item Why can't you change a streaming aggregation's schema freely on restart?
    \item Which metric tells you the stream is falling behind?
\end{itemize}}

\prodnote{One checkpoint location, one stream definition: pointing a second query at the same checkpoint corrupts offsets. Treat checkpoint paths as part of the pipeline contract --- version them in your repo README and govern them in Unity Catalog volumes.}

\subsection*{Week 14: State Management, Watermarks \& Exactly-Once}
\par\noindent\textbf{Outcomes:} implement event-time windows with watermarks, manage stateful aggregations and deduplication (\texttt{dropDuplicates}, \texttt{withWatermark}), use \texttt{flatMapGroupsWithState} and the RocksDB state store, and reason about exactly-once end-to-end.
\begin{itemize}[label=-]
    \dayitem{Monday}{Event time vs.\ processing time; tumbling, sliding, and session windows.}
    \dayitem{Tuesday}{Watermarks: late-data thresholds, state eviction, and their interaction with output modes.}
    \dayitem{Wednesday}{Stateful operators: deduplication, arbitrary stateful processing, state store backends (HDFS vs.\ RocksDB), and exactly-once composition with idempotent sinks.}
    \dayitem{Thursday}{\project{Build a windowed aggregation over out-of-order IoT events with a 10-minute watermark, inject late and duplicate events, prove dedup via \texttt{dropDuplicates} on \texttt{event\_id}, and compare state-store memory with and without RocksDB.}}
    \dayitem{Friday}{\usecase{Design a fraud-detection stream that must tolerate 15 minutes of late events, deduplicate replays, and still emit alerts within 30 seconds of event time.}}
\end{itemize}
\resources{Structured Streaming watermarking docs (\url{https://spark.apache.org/docs/latest/structured-streaming-programming-guide.html\#handling-late-data-and-watermarking}); RocksDB state store guide (\url{https://docs.databricks.com/structured-streaming/rocksdb-state-store.html}).}
\selfcheck{\begin{itemize}[label=-, leftmargin=1.5em, itemsep=0pt]
    \item What does a watermark actually delete from memory?
    \item Why does exactly-once require both a checkpoint \emph{and} an idempotent sink?
    \item Tumbling vs.\ sliding window: what does the slide add?
    \item When must you switch from the default state store to RocksDB?
    \item Why does \texttt{dropDuplicates} without a watermark risk unbounded state?
\end{itemize}}

\subsection*{Week 15: Kafka, Event Hubs \& Kinesis Integration}
\par\noindent\textbf{Outcomes:} connect Structured Streaming to Apache Kafka (and Kafka-compatible Event Hubs) and Amazon Kinesis, manage offsets and \texttt{startingOffsets}/\texttt{maxOffsetsPerTrigger}, deserialize JSON/Avro, and design consumer groups for scale.
\begin{itemize}[label=-]
    \dayitem{Monday}{Kafka architecture for data engineers: topics, partitions, consumer groups, and offset semantics.}
    \dayitem{Tuesday}{Spark Kafka source: options, SSL/SASL authentication, schema registry and Avro, Event Hubs via the Kafka endpoint.}
    \dayitem{Wednesday}{Kinesis source, rate limiting with \texttt{maxOffsetsPerTrigger}, backpressure, and multi-stream topologies.}
    \dayitem{Thursday}{\project{Run a local Kafka (Docker) producing synthetic transactions, consume with Structured Streaming into Bronze Delta, enforce a 1k-records-per-trigger cap, then replay from \texttt{startingOffsets="earliest"} into a recovery table and reconcile counts.}}
    \dayitem{Friday}{\usecase{Design the messaging backbone for a retailer migrating from on-prem Kafka to Event Hubs while keeping Databricks consumers unchanged and zero-downtime.}}
\end{itemize}
\resources{Spark Kafka integration guide (\url{https://spark.apache.org/docs/latest/structured-streaming-kafka-integration.html}); Event Hubs Kafka endpoint docs (\url{https://learn.microsoft.com/azure/event-hubs/event-hubs-for-kafka-ecosystem-overview}).}
\selfcheck{\begin{itemize}[label=-, leftmargin=1.5em, itemsep=0pt]
    \item Why does Kafka give at-least-once delivery, and where do you add idempotence?
    \item What does \texttt{maxOffsetsPerTrigger} protect?
    \item Why prefer a schema registry over inferred JSON schemas for streaming?
    \item How do Kafka partitions map to Spark tasks?
    \item What is the Event Hubs trick that lets Kafka clients connect unchanged?
\end{itemize}}

\subsection*{Week 16: Real-Time Lakehouse Architectures}
\par\noindent\textbf{Outcomes:} design streaming medallion architectures, combine stream and batch in one Delta table, serve real-time dashboards from Gold streaming tables, and operate latency/cost trade-offs deliberately.
\begin{itemize}[label=-]
    \dayitem{Monday}{Streaming medallion: Bronze append streams, Silver streaming joins and dedup, Gold incremental aggregates.}
    \dayitem{Tuesday}{Stream-static and stream-stream joins, CDC-to-stream fan-in, and \texttt{foreachBatch} for merge-based sinks.}
    \dayitem{Wednesday}{Serving real-time Gold to AI/BI dashboards, latency budgets, and cost models for continuous vs.\ triggered processing.}
    \dayitem{Thursday}{\project{Build a real-time orders pipeline: Kafka $\rightarrow$ Bronze stream $\rightarrow$ Silver deduped stream with \texttt{foreachBatch} MERGE $\rightarrow$ Gold 5-minute tumbling revenue table, visualized on an AI/BI dashboard refreshing every minute.}}
    \dayitem{Friday}{\usecase{Architect a real-time lakehouse for a logistics fleet: 10k vehicles $\times$ 1 event/second, sub-minute dispatch dashboards, and a daily batch reconciliation that must agree with the stream.}}
\end{itemize}
\resources{Databricks real-time architecture patterns (\url{https://www.databricks.com/blog}); Structured Streaming + Delta guide (\url{https://docs.databricks.com/structured-streaming/delta-lake.html}).}
\selfcheck{\begin{itemize}[label=-, leftmargin=1.5em, itemsep=0pt]
    \item Why is \texttt{foreachBatch} the escape hatch for MERGE sinks?
    \item How can one Delta table serve both stream and batch writers safely?
    \item What latency budget forces continuous processing over 1-minute triggers?
    \item Why keep a batch reconciliation path even in a ``real-time'' system?
    \item What does a streaming Gold table cost per day vs.\ a triggered one?
\end{itemize}}

\capstone{Milestone 4}{Build a Real-Time Fraud Telemetry Platform. Stream synthetic card transactions from Kafka into a Bronze Delta table, deduplicate and window in Silver with a watermark, flag impossible-travel fraud with a stateful rule, publish Gold 5-minute aggregates to a live AI/BI dashboard, and prove exactly-once via a replay drill with zero duplicate alerts.}
\rubric{\begin{itemize}[label=-, leftmargin=1.5em, itemsep=0pt]
    \item Replay drill: re-consuming the Kafka topic produces zero duplicate Silver rows and zero duplicate alerts.
    \item Watermark + dedup configured and demonstrated with injected late/duplicate events.
    \item Impossible-travel rule flags a planted fraud case within the latency budget.
    \item Dashboard refreshes at decision cadence (not platform maximum) with a freshness indicator.
    \item Cost estimate: continuous vs.\ triggered DBU comparison with the chosen trigger justified.
    \item Security review: Kafka credentials in secret scopes, checkpoint governance, no PII in Bronze beyond necessity.
\end{itemize}}

\subsection*{Milestone 4 Capstone: Reference Solution Sketch}
\begin{itemize}[label=-, itemsep=0pt]
    \item \textbf{Architecture:} Kafka $\rightarrow$ Bronze (append stream) $\rightarrow$ Silver (\texttt{dropDuplicates} + watermark + \texttt{foreachBatch} MERGE) $\rightarrow$ Gold tumbling aggregates $\rightarrow$ AI/BI dashboard; the stateful impossible-travel rule runs as a second Silver stream.
    \item \textbf{Steps:} stand up the Kafka producer with unique \texttt{event\_id}s; configure checkpoints per query; set \texttt{maxOffsetsPerTrigger}; implement dedup before rules; build the dashboard; run the replay drill.
    \item \textbf{Hardest part --- dedup-before-rules ordering:}
\begin{lstlisting}
silver = (bronze_stream
  .withWatermark("event_ts", "10 minutes")
  .dropDuplicates(["event_id"])
  .groupBy(window("event_ts", "5 minutes"), "card_id")
  .agg(count("*").alias("txns"),
       countDistinct("city").alias("cities")))
\end{lstlisting}
    \item \textbf{Pitfalls:} rules reading the raw (undeduplicated) stream double-fire on replay; watermark shorter than real lateness silently drops valid events; a continuous trigger left running over the weekend burns the week's budget.
\end{itemize}

\begin{figure}[htbp]
\centering
\begin{tikzpicture}[
    node distance=0.55cm and 0.9cm,
    box/.style={draw, rounded corners, fill=blue!10, align=center, font=\small, minimum height=0.9cm, inner sep=4pt},
    gate/.style={draw, rounded corners, fill=orange!15, align=center, font=\footnotesize, inner sep=3pt},
    arr/.style={-{Stealth}, thick}
]
\node[box] (kafka) {Kafka\\transactions};
\node[box, right=of kafka, fill=brown!12] (bronze) {Bronze\\append stream};
\node[box, right=of bronze, fill=gray!15] (silver) {Silver\\dedup + watermark};
\node[gate, below=0.5cm of silver] (rule) {Impossible-travel\\stateful rule};
\node[box, right=of silver, fill=yellow!18] (gold) {Gold\\5-min aggregates};
\node[box, right=of gold] (dash) {AI/BI\\dashboard};
\draw[arr] (kafka) -- (bronze);
\draw[arr] (bronze) -- (silver);
\draw[arr] (silver) -- (gold);
\draw[arr] (gold) -- (dash);
\draw[arr, dashed] (silver) -- (rule);
\end{tikzpicture}
\caption{Milestone 4 reference architecture: real-time fraud telemetry lakehouse}
\end{figure}

% ==========================================
% MILESTONE 5
% ==========================================
\section{Milestone 5: ML \& MLOps (Month 5)}
\textit{Objective: Productionize machine learning on Databricks with MLflow, the Feature Store, Mosaic AI, and model serving.}
\par\noindent\textbf{Learning outcomes:} track experiments and register models with MLflow, build leakage-safe feature tables, assemble a RAG pipeline with Vector Search, and serve and monitor models with governance.

\subsection*{Week 17: MLflow \& Experiment Tracking}
\par\noindent\textbf{Outcomes:} track parameters, metrics, and artifacts with MLflow autologging, compare runs in the UI, register models in Unity Catalog, and reproduce a training run from its logged environment.
\begin{itemize}[label=-]
    \dayitem{Monday}{MLflow components: Tracking, Models, Registry (in Unity Catalog), and the run/artifact model.}
    \dayitem{Tuesday}{Autologging, manual logging, run tagging with dataset version and git commit, and model signatures.}
    \dayitem{Wednesday}{Model registry workflows: aliases (champion/challenger), versioning, approval, and loading models by alias.}
    \dayitem{Thursday}{\project{Train six scikit-learn churn models with MLflow autologging, tag each run with dataset snapshot and commit hash, register the best as \texttt{champion}, and load it back in a fresh notebook via \texttt{models:/catalog.schema.churn\_model@champion}.}}
    \dayitem{Friday}{\usecase{Design the experiment-governance process for a bank where every production model must trace to data, code, hyperparameters, and an approver.}}
\end{itemize}
\resources{MLflow documentation (\url{https://mlflow.org/docs/latest/}); MLflow on Databricks guide (\url{https://docs.databricks.com/mlflow/}).}
\selfcheck{\begin{itemize}[label=-, leftmargin=1.5em, itemsep=0pt]
    \item What does autologging capture that manual logging usually forgets?
    \item Why register models in Unity Catalog instead of the workspace registry?
    \item Alias vs.\ version number: which should production jobs reference?
    \item What makes a run reproducible a year later?
    \item Why compare runs only on identical data snapshots?
\end{itemize}}

\subsection*{Week 18: Feature Engineering \& Feature Store}
\par\noindent\textbf{Outcomes:} build feature tables with \texttt{FeatureEngineeringClient}, enforce point-in-time correctness with \texttt{create\_training\_set}, serve online features, and prevent training-serving skew.
\begin{itemize}[label=-]
    \dayitem{Monday}{Feature engineering discipline: leakage, point-in-time correctness, and feature definitions as code.}
    \dayitem{Tuesday}{Feature Store in Unity Catalog: feature tables, \texttt{create\_training\_set} with timestamp lookup, and lineage to models.}
    \dayitem{Wednesday}{Online stores and feature serving, on-demand features, and feature freshness monitoring.}
    \dayitem{Thursday}{\project{Build \texttt{customer\_features} and \texttt{order\_features} tables with as-of timestamps, create a training set with point-in-time joins, train a model that records feature-table lineage, and verify the same features are served at inference.}}
    \dayitem{Friday}{\usecase{Audit a churn model that performs 15 points better in training than production: find the leakage and redesign the feature pipeline.}}
\end{itemize}
\resources{Feature Engineering in Unity Catalog docs (\url{https://docs.databricks.com/machine-learning/feature-store/}); point-in-time join guide (\url{https://docs.databricks.com/machine-learning/feature-store/time-series.html}).}
\selfcheck{\begin{itemize}[label=-, leftmargin=1.5em, itemsep=0pt]
    \item What is label leakage and give one concrete example?
    \item How does a timestamp lookup prevent future data entering training rows?
    \item Why does the registry record feature-table lineage?
    \item Online vs.\ offline feature store: which serves inference at 10 ms?
    \item What is training-serving skew and how does the Feature Store prevent it?
\end{itemize}}

\subsection*{Week 19: Mosaic AI, Vector Search \& GenAI (RAG)}
\par\noindent\textbf{Outcomes:} build embeddings pipelines, create and query Mosaic AI Vector Search indexes, assemble a retrieval-augmented generation (RAG) chain, and evaluate answer quality with MLflow evaluation.
\begin{itemize}[label=-]
    \dayitem{Monday}{GenAI building blocks on Databricks: Foundation Model APIs, embeddings, and Model Serving endpoints.}
    \dayitem{Tuesday}{Vector Search: managed vs.\ Delta-sync indexes, embedding models, similarity search, and metadata filters.}
    \dayitem{Wednesday}{RAG architecture: chunking, retrieval, prompt assembly, evaluation with MLflow, and governance of GenAI pipelines.}
    \dayitem{Thursday}{\project{Build a RAG assistant over this program's own documentation: chunk markdown into a Delta table, sync a Vector Search index, retrieve top-k context for a question, and assemble an answer via a Foundation Model endpoint; evaluate 10 questions with MLflow and log the quality metrics.}}
    \dayitem{Friday}{\usecase{Design a customer-support copilot over 2M help-center articles with freshness requirements, PII redaction, and a cost ceiling per answer.}}
\end{itemize}
\resources{Mosaic AI documentation (\url{https://docs.databricks.com/machine-learning/}); Vector Search documentation (\url{https://docs.databricks.com/generative-ai/vector-search.html}).}
\selfcheck{\begin{itemize}[label=-, leftmargin=1.5em, itemsep=0pt]
    \item Why store source chunks in Delta and sync the index instead of writing vectors directly?
    \item Managed vs.\ Delta-sync index: which do you choose for a table that updates hourly?
    \item What does chunk size trade off in retrieval quality?
    \item Why evaluate RAG with logged metrics rather than eyeballing answers?
    \item Where should PII redaction happen in the RAG pipeline?
\end{itemize}}

\subsection*{Week 20: Model Serving, Monitoring \& MLOps}
\par\noindent\textbf{Outcomes:} deploy models to serverless serving endpoints, implement batch scoring with idempotent writes, monitor input drift and prediction quality with Lakehouse Monitoring, and gate promotion behind metric thresholds.
\begin{itemize}[label=-]
    \dayitem{Monday}{Model Serving: serverless endpoints, scaling to zero, provisioned concurrency, and endpoint governance.}
    \dayitem{Tuesday}{Batch and streaming scoring patterns: alias-loaded models as Spark UDFs, partition-idempotent writes.}
    \dayitem{Wednesday}{Lakehouse Monitoring: inference tables, drift metrics, alerting, and the retraining loop with human-gated promotion.}
    \dayitem{Thursday}{\project{Deploy the Week 17 champion to a serverless endpoint, run a batch-scoring job over 100k rows with \texttt{replaceWhere} idempotency, enable an inference table, inject skewed inputs, and prove a monitor catches the drift while a retraining trigger fires but does not auto-promote.}}
    \dayitem{Friday}{\usecase{Design the MLOps control plane for a lender whose regulator requires drift evidence, approval records, and rollback for every model in production.}}
\end{itemize}
\resources{Model Serving documentation (\url{https://docs.databricks.com/machine-learning/model-serving/}); Lakehouse Monitoring documentation (\url{https://docs.databricks.com/lakehouse-monitoring/}).}
\selfcheck{\begin{itemize}[label=-, leftmargin=1.5em, itemsep=0pt]
    \item What does an inference table capture and why does that enable monitoring?
    \item Why load models by alias in scoring jobs?
    \item What makes a batch-scoring write idempotent?
    \item Why must retraining triggers not auto-promote to production?
    \item Scale-to-zero vs.\ provisioned concurrency: which fits a latency-critical endpoint?
\end{itemize}}

\capstone{Milestone 5}{Build an End-to-End MLOps Pipeline. Engineer leakage-safe features into Unity Catalog feature tables, train and track churn models with MLflow, register a champion, build a RAG knowledge assistant over the platform's runbooks with Vector Search, serve the model on a serverless endpoint, and operate drift monitoring with a human-gated retraining loop.}
\rubric{\begin{itemize}[label=-, leftmargin=1.5em, itemsep=0pt]
    \item Training set demonstrably point-in-time correct; one injected leakage feature caught by review or test.
    \item All runs tagged with dataset snapshot + commit; champion loads by alias in a clean notebook.
    \item RAG pipeline answers 10 evaluation questions with logged MLflow quality metrics.
    \item Batch scoring is idempotent (\texttt{replaceWhere}); endpoint serves a live request with latency recorded.
    \item Drift monitor fires on injected skew; retraining triggered but promotion blocked pending human approval.
    \item Security review: endpoint access controlled, secrets in scopes, PII handling in features documented.
\end{itemize}}

\subsection*{Milestone 5 Capstone: Reference Solution Sketch}
\begin{itemize}[label=-, itemsep=0pt]
    \item \textbf{Architecture:} Gold $\rightarrow$ feature tables (as-of timestamps) $\rightarrow$ MLflow-tracked training $\rightarrow$ UC registry (champion alias) $\rightarrow$ serverless endpoint + inference table $\rightarrow$ Lakehouse monitor $\rightarrow$ gated retraining; RAG branch: docs $\rightarrow$ chunks Delta $\rightarrow$ Vector Search $\rightarrow$ Foundation Model endpoint.
    \item \textbf{Steps:} author feature jobs with \texttt{as\_of\_date} parameters; build training set with timestamp lookups; autolog six runs; promote by alias; deploy endpoint; enable monitoring; inject drift; observe gate.
    \item \textbf{Hardest part --- point-in-time training set:}
\begin{lstlisting}
training_set = fs.create_training_set(
    df=labels_df,               # has customer_id + event_ts
    feature_lookups=[
        FeatureLookup("ml.customer_features",
                      lookup_key="customer_id",
                      timestamp_lookup_key="event_ts")],
    label="churned",
    exclude_columns=["customer_id", "event_ts"])
\end{lstlisting}
    \item \textbf{Pitfalls:} features computed past their as-of date (leakage by construction); scoring against ad hoc features instead of the governed table; an unmonitored endpoint whose drift you discover from a business stakeholder.
\end{itemize}

\begin{figure}[htbp]
\centering
\begin{tikzpicture}[
    node distance=0.55cm and 0.85cm,
    box/.style={draw, rounded corners, fill=blue!10, align=center, font=\small, minimum height=0.9cm, inner sep=4pt},
    ml/.style={draw, rounded corners, fill=magenta!10, align=center, font=\footnotesize, inner sep=3pt},
    arr/.style={-{Stealth}, thick}
]
\node[box] (gold) {Gold tables};
\node[box, right=of gold] (feat) {Feature tables\\(as-of joins)};
\node[ml, right=of feat] (mlflow) {MLflow runs\\+ UC registry};
\node[box, right=of mlflow] (serve) {Serving endpoint\\+ inference table};
\node[ml, below=0.5cm of serve] (mon) {Lakehouse Monitoring\\drift alarms};
\node[ml, above=0.5cm of mlflow] (rag) {Vector Search\\RAG assistant};
\draw[arr] (gold) -- (feat);
\draw[arr] (feat) -- (mlflow);
\draw[arr] (mlflow) -- (serve);
\draw[arr] (serve) -- (mon);
\draw[arr, dashed] (mon.west) -| (mlflow.south);
\draw[arr, dashed] (gold) -- (rag);
\end{tikzpicture}
\caption{Milestone 5 reference architecture: governed MLOps with monitoring and RAG}
\end{figure}

% ==========================================
% MILESTONE 6
% ==========================================
\section{Milestone 6: Security, Production \& Certification (Month 6)}
\textit{Objective: Harden the platform, operate deep governance, automate delivery with Asset Bundles and Terraform, and prepare for certification.}
\par\noindent\textbf{Learning outcomes:} implement identity and network isolation, exploit Unity Catalog lineage/audit/sharing, deploy with Databricks Asset Bundles and Terraform, optimize cost, and pass the Databricks Certified Data Engineer exam.

\subsection*{Week 21: Security, IAM, SCIM \& Network Isolation}
\par\noindent\textbf{Outcomes:} design account/workspace identity with SSO and SCIM provisioning, use service principals and secret scopes, configure network isolation (Private Link/PrivateLink-style connectivity, IP access lists), and apply least privilege across compute and data.
\begin{itemize}[label=-]
    \dayitem{Monday}{Identity architecture: account console vs.\ workspace, SSO with the IdP, SCIM provisioning, and group strategy.}
    \dayitem{Tuesday}{Service principals, OAuth machine-to-machine, personal access token policy, and secret scopes (Databricks-backed vs.\ Key Vault-backed).}
    \dayitem{Wednesday}{Network security: private connectivity, IP access lists, storage firewall rules, and audit-ready network design.}
    \dayitem{Thursday}{\project{Provision a service principal via the CLI, grant it only the schemas it needs, store its credentials in a secret scope, convert a notebook from a hardcoded PAT to \texttt{dbutils.secrets.get}, and add an IP access list; verify a blocked connection is denied.}}
    \dayitem{Friday}{\usecase{Design the security architecture for a multi-team platform where contractors get 90-day access, production jobs run as service principals, and all traffic must stay off the public internet.}}
\end{itemize}
\resources{Databricks security documentation (\url{https://docs.databricks.com/security/}); SCIM provisioning guide (\url{https://docs.databricks.com/admin/users-groups/scim/}).}
\selfcheck{\begin{itemize}[label=-, leftmargin=1.5em, itemsep=0pt]
    \item Why do production jobs run as service principals, not people?
    \item What does SCIM automate that manual user management cannot?
    \item Databricks-backed vs.\ Key Vault-backed secret scope: who controls rotation?
    \item What does an IP access list \emph{not} protect against?
    \item How does private connectivity change your threat model?
\end{itemize}}

\subsection*{Week 22: Unity Catalog Deep Dive: Lineage, Audit \& Delta Sharing}
\par\noindent\textbf{Outcomes:} read column-level lineage from system tables, build an audit practice on \texttt{system.access.audit}, share data securely with Delta Sharing (open protocol and Databricks-to-Databricks), and implement data classification and tag-based policies.
\begin{itemize}[label=-]
    \dayitem{Monday}{Lineage: table and column lineage, capturing it, and using it for impact analysis and debugging.}
    \dayitem{Tuesday}{Audit deep dive: \texttt{system.access.audit} schema, alerting on sensitive access, and compliance reporting.}
    \dayitem{Wednesday}{Delta Sharing: shares, recipients, clean rooms, and open-protocol vs.\ Databricks-to-Databricks sharing.}
    \dayitem{Thursday}{\project{Build a lineage report tracing \texttt{gold.fact\_orders.total\_amount} back to Bronze columns, write an audit alert for any read of the PII-mapped columns by a non-approved group, and share a Gold table to an external recipient via Delta Sharing with a documented revocation test.}}
    \dayitem{Friday}{\usecase{Design data sharing for a supply-chain consortium: five companies exchange curated datasets with auditability, revocation, and no cross-company credentials.}}
\end{itemize}
\resources{Unity Catalog lineage documentation (\url{https://docs.databricks.com/data-governance/unity-catalog/data-lineage.html}); Delta Sharing documentation (\url{https://docs.databricks.com/delta-sharing/}).}
\selfcheck{\begin{itemize}[label=-, leftmargin=1.5em, itemsep=0pt]
    \item How does column-level lineage accelerate a breaking-change impact analysis?
    \item Which system table records ``who ran what query''?
    \item Why does Delta Sharing avoid giving consumers your storage credentials?
    \item What happens to a recipient's access when you revoke a share?
    \item Why is open-protocol sharing important for non-Databricks consumers?
\end{itemize}}

\subsection*{Week 23: CI/CD with Databricks Asset Bundles, Terraform \& Cost Optimization}
\par\noindent\textbf{Outcomes:} define jobs/pipelines as code with Databricks Asset Bundles (YAML), provision infrastructure with the Terraform Databricks provider, wire GitHub Actions CI/CD with dev/staging/prod targets, and build a cost-optimization practice from system billing tables.
\begin{itemize}[label=-]
    \dayitem{Monday}{Databricks Asset Bundles: \texttt{databricks.yml}, resources, targets, variables, and \texttt{bundle deploy/validate/run}.}
    \dayitem{Tuesday}{Terraform provider for Databricks: workspace resources, drift detection, and state management.}
    \dayitem{Wednesday}{CI/CD pipelines: GitHub Actions with OIDC/service-principal auth, PR validation, and dev$\rightarrow$staging$\rightarrow$prod promotion; cost analysis with \texttt{system.billing.usage}.}
    \dayitem{Thursday}{\project{Convert the Month 3 Workflow into an Asset Bundle with dev and prod targets, validate and deploy from GitHub Actions on merge, provision the target workspace resources with Terraform, and produce a weekly DBU cost report with two implemented optimizations.}}
    \dayitem{Friday}{\usecase{Design the delivery and FinOps operating model for a platform team shipping 30 pipelines across three environments with monthly budget accountability per domain.}}
\end{itemize}
\resources{Databricks Asset Bundles documentation (\url{https://docs.databricks.com/dev-tools/bundles/}); Terraform Databricks provider (\url{https://registry.terraform.io/providers/databricks/databricks/latest/docs}).}
\selfcheck{\begin{itemize}[label=-, leftmargin=1.5em, itemsep=0pt]
    \item What does \texttt{databricks bundle validate} catch before deployment?
    \item Why separate bundle targets per environment rather than parameter flags?
    \item Terraform state: why must it be remote and locked?
    \item Which system table powers chargeback reports?
    \item Name two compute settings that cut DBU spend without slowing jobs.
\end{itemize}}

\prodnote{Pin the deployed branch, not a floating one: a bundle that deploys \texttt{main} on every merge without a staging gate reproduces the classic ``Friday 5 p.m.\ broken prod pipeline.'' Require the staging target to run green in CI before prod deployment is permitted.}

\subsection*{Week 24: Final Capstone \& Databricks Data Engineer Certification Prep}
\par\noindent\textbf{Outcomes:} integrate all six milestones into one production-grade platform, execute a certification study plan for the Databricks Certified Data Engineer Associate/Professional, and complete two timed mock exams.
\begin{itemize}[label=-]
    \dayitem{Monday}{Full-platform review: lakehouse, warehousing, ETL, streaming, ML, and governance --- one integrated architecture map.}
    \dayitem{Tuesday}{Certification domains walkthrough: Databricks platform, Delta Lake, ETL with Spark/DLT, production pipelines, and governance.}
    \dayitem{Wednesday}{Exam drills: scenario questions, time management, and weak-domain remediation by rebuilding labs.}
    \dayitem{Thursday}{\project{Final Capstone build (see below): the complete governed data platform deployed from an Asset Bundle with CI/CD, monitoring, cost report, and security review.}}
    \dayitem{Friday}{\usecase{Present the final capstone as a 15-minute architecture review defending three design decisions with trade-offs, one failure story, and cost numbers.}}
\end{itemize}
\resources{Databricks certification page (\url{https://www.databricks.com/learn/certification}); Databricks Academy self-paced courses (\url{https://www.databricks.com/learn/training}).}
\selfcheck{\begin{itemize}[label=-, leftmargin=1.5em, itemsep=0pt]
    \item Which exam domain is your weakest, and which lab rebuilds it?
    \item Can you deploy your whole platform from a clean clone? Prove it.
    \item What three numbers (cost, latency, rows) summarize your platform?
    \item Which single design decision would you revisit at 10x scale?
    \item Can you explain every box on your architecture diagram in one sentence?
\end{itemize}}

\capstone{Milestone 6}{Final Capstone: The Governed Data Platform. Deploy a complete platform --- medallion lakehouse with optimized Delta storage, DLT pipelines with expectations, Auto Loader + Kafka streaming, MLflow-managed model with serving and monitoring, Unity Catalog governance with lineage/audit/Delta Sharing --- as code via Asset Bundles and Terraform, promoted through dev/staging/prod by GitHub Actions, with a cost report and a certification-style architecture defense.}
\rubric{\begin{itemize}[label=-, leftmargin=1.5em, itemsep=0pt]
    \item Clean-environment redeploy: \texttt{terraform apply} + \texttt{databricks bundle deploy -t prod} rebuilds the platform from the repo alone.
    \item Batch and streaming paths both feed Gold with reconciliation; replay drill produces zero duplicates.
    \item Governance evidence: access matrix, lineage report, audit alert, and one Delta Sharing share with revocation test.
    \item CI/CD: PR validation, staging gate, prod promotion on merge; rollback documented.
    \item Cost report from \texttt{system.billing.usage} with per-domain chargeback and two implemented optimizations.
    \item Two consecutive 80\%+ mock exam scores with no domain below 70\%.
\end{itemize}}

\subsection*{Milestone 6 Capstone: Reference Solution Sketch}
\begin{itemize}[label=-, itemsep=0pt]
    \item \textbf{Architecture:} one bundle defining volumes, pipelines, jobs, serving endpoint, and monitors across three targets; Terraform provisioning workspace-level resources; Unity Catalog enforcing the access matrix; system tables powering audit and cost reporting.
    \item \textbf{Steps:} extract every hand-built artifact from Months 1--5 into bundle YAML; define \texttt{dev}/\texttt{staging}/\texttt{prod} targets with variables; wire GitHub Actions with service-principal auth; run the redeploy drill in a fresh workspace; present.
    \item \textbf{Hardest part --- the bundle skeleton:}
\begin{lstlisting}
bundle:
  name: governed-platform
targets:
  prod:
    mode: production
    workspace:
      host: https://<workspace>.azuredatabricks.net
resources:
  jobs:
    nightly_elt:
      name: nightly-elt
      tasks:
        - task_key: ingest
          pipeline_task: {pipeline_id: \${resources.pipelines.medallion.id}}
\end{lstlisting}
    \item \textbf{Pitfalls:} hardcoded workspace IDs breaking the prod target; Terraform state in a local file; a demo that works only because your personal credentials are cached --- test with the service principal.
\end{itemize}

\begin{figure}[htbp]
\centering
\begin{tikzpicture}[
    node distance=0.55cm and 0.85cm,
    box/.style={draw, rounded corners, fill=blue!10, align=center, font=\small, minimum height=0.9cm, inner sep=4pt},
    ops/.style={draw, rounded corners, fill=teal!10, align=center, font=\footnotesize, inner sep=3pt},
    arr/.style={-{Stealth}, thick}
]
\node[box] (src) {Batch files\\+ Kafka};
\node[box, right=of src] (lake) {Medallion\\lakehouse (DLT)};
\node[box, right=of lake] (serve) {SQL warehouse\\+ ML serving};
\node[ops, below=0.5cm of lake] (uc) {Unity Catalog: lineage, audit, sharing};
\node[ops, above=0.5cm of lake] (cicd) {Asset Bundles + Terraform\\GitHub Actions CI/CD};
\draw[arr] (src) -- (lake);
\draw[arr] (lake) -- (serve);
\draw[arr, dashed] (uc) -- (lake);
\draw[arr, dashed] (cicd) -- (lake);
\end{tikzpicture}
\caption{Milestone 6 reference architecture: the complete governed, code-deployed platform}
\end{figure}

% ==========================================
\section{Appendix A: Glossary}
\begin{description}[leftmargin=3.5cm, style=nextline, itemsep=1pt]
    \item[Delta Lake] Open table format adding ACID transactions, time travel, and schema enforcement to Parquet files in cloud storage.
    \item[Transaction Log] The \texttt{\_delta\_log} directory of JSON commits (plus checkpoints) making Delta writes atomic and versions queryable.
    \item[Medallion Architecture] Progressive-refinement lakehouse pattern: Bronze (raw), Silver (validated), Gold (curated).
    \item[Unity Catalog] Databricks' unified governance layer: metastore $\rightarrow$ catalog $\rightarrow$ schema $\rightarrow$ table/volume, with lineage and audit.
    \item[Volume] Unity Catalog object governing access to files in cloud storage for non-tabular data.
    \item[DBU] Databricks Unit --- the compute billing unit consumed per second of cluster or serverless runtime.
    \item[Photon] Databricks' vectorized C++ query engine accelerating SQL and DataFrame workloads.
    \item[SQL Warehouse] SQL-optimized compute (serverless, pro, or classic) powering Databricks SQL, dashboards, and Genie.
    \item[AI/BI Dashboard] Databricks-native BI artifact with parameterized datasets and visualizations on governed data.
    \item[Genie] AI/BI natural-language interface that generates SQL over curated, instructed datasets.
    \item[OPTIMIZE] Delta command compacting small files, optionally with \texttt{ZORDER BY} multi-dimensional clustering.
    \item[Liquid Clustering] Delta layout technique (\texttt{CLUSTER BY}) replacing fixed partitions/Z-ORDER with rewritable-free clustering keys.
    \item[Deletion Vector] Delta row-level delete marker avoiding full file rewrites for \texttt{DELETE}/\texttt{UPDATE}/\texttt{MERGE}.
    \item[Auto Loader] Incremental file-ingestion source (\texttt{cloudFiles}) with checkpointed exactly-once semantics and schema evolution.
    \item[DLT (Delta Live Tables)] Declarative pipeline framework with streaming tables, materialized views, and data-quality expectations.
    \item[Expectation] DLT data-quality rule with warn (\texttt{expect}), drop, or fail semantics.
    \item[Checkpoint] Structured Streaming (or Auto Loader) state directory enabling exactly-once recovery after restarts.
    \item[Watermark] Event-time threshold bounding how late events are accepted and when state is evicted.
    \item[Task Value] Runtime value passed between tasks in a Databricks Workflow job.
    \item[MLflow] Open-source ML lifecycle platform: tracking, models, registry (in Unity Catalog on Databricks), and evaluation.
    \item[Feature Store] Governed feature tables with point-in-time-correct training sets and serving lineage.
    \item[Mosaic AI Vector Search] Managed vector index over Delta tables powering semantic retrieval for RAG.
    \item[Model Serving] Serverless (or provisioned) endpoints exposing registered models, with inference tables for monitoring.
    \item[Lakehouse Monitoring] Managed drift and quality monitoring over Delta tables and model inference logs.
    \item[Delta Sharing] Open protocol for sharing live Delta tables to any recipient without copying data or sharing storage credentials.
    \item[Asset Bundle] Databricks' infrastructure-as-code format (\texttt{databricks.yml}) defining jobs, pipelines, and targets for CI/CD.
    \item[Service Principal] Non-human identity for automated jobs and CI/CD, granted least-privilege access.
    \item[Secret Scope] Governed vault reference (Databricks- or Key Vault-backed) injected via \texttt{dbutils.secrets.get}.
\end{description}

\section{Appendix B: Assessment \& Grading Model}
Each monthly capstone is graded on the following weighted rubric:

\begin{center}
\renewcommand{\arraystretch}{1.3}
\begin{tabular}{@{} l c p{7.5cm} @{}}
\toprule
\textbf{Category} & \textbf{Weight} & \textbf{What is evaluated} \\
\midrule
Correctness & 30\% & Pipeline runs end-to-end and produces the specified outputs from a clean environment. \\
Reliability \& tests & 20\% & Data-quality expectations, failure handling, retries/repair runs, and a demonstrated failing case. \\
Security & 15\% & Least-privilege service principals, secret scopes, no hardcoded tokens, PII masking. \\
Cost awareness & 15\% & DBU-based monthly estimate with at least 2 documented optimization decisions. \\
Documentation & 10\% & README, architecture diagram, and runbook quality. \\
Demo & 10\% & Live or recorded walkthrough of the working system. \\
\bottomrule
\end{tabular}
\end{center}

The two full-length mock exams scheduled for the weekend after Week 24 feed the final certification readiness score: a consistent score of 80\%+ across both mocks indicates readiness to book the Databricks Certified Data Engineer exam; below that, revisit the weak domains noted during Week 24 Friday before booking.

\section{Appendix C: Interview Preparation}
\begin{enumerate}[itemsep=1pt]
    \item Explain the Delta transaction log and why it makes a data lake ACID-compliant. (Week 2)
    \item When does liquid clustering beat partition + Z-ORDER, and when would you keep partitioning? (Week 5)
    \item A Databricks SQL dashboard burst at 9 a.m.\ queues for minutes --- how do you fix it? (Week 6)
    \item Design Unity Catalog grants so analysts read masked Gold while engineers write Silver. (Week 8)
    \item A PySpark join is skewed on one key --- walk through your Spark UI evidence and the salting fix. (Week 9)
    \item DLT vs.\ plain Spark + Workflows: where do you draw the line, and how do expectations change your testing? (Weeks 10, 12)
    \item How does Auto Loader guarantee exactly-once ingestion, and what breaks if someone deletes the checkpoint? (Week 11)
    \item Design a streaming fraud pipeline with watermarks, dedup, and exactly-once semantics --- end to end. (Weeks 13--14)
    \item How do you prove a training set has no label leakage, and what does the Feature Store contribute? (Week 18)
    \item Design a governed multi-environment delivery model with Asset Bundles, Terraform, and per-domain chargeback. (Weeks 22--23, Capstone 6)
\end{enumerate}

\newpage
% ==========================================
\section{Appendix D: Self-Check Answer Key}
Terse answers to the weekly self-check questions.

\noindent\textbf{Week 1:}
\begin{enumerate}[itemsep=0pt, leftmargin=2em]
    \item Workspace metadata, notebooks, and the web UI live in the control plane; cluster VMs and your data live in the data plane.
    \item Job clusters start and stop per run and are priced lower; all-purpose clusters bill for interactive uptime.
    \item It caps configuration choices (size, auto-termination, spot use), protecting budgets and standards.
    \item A Git folder syncs with a remote Git branch; a plain workspace folder is not version-controlled.
    \item OAuth issues short-lived tokens; a PAT is a long-lived secret that can leak from scripts.
\end{enumerate}
\noindent\textbf{Week 2:}
\begin{enumerate}[itemsep=0pt, leftmargin=2em]
    \item The commit to \texttt{\_delta\_log} is a single atomic file write listing all data files for the version.
    \item Optimistic concurrency with snapshot isolation; conflicting concurrent writes retry or fail.
    \item \texttt{mergeSchema} for additive evolution; \texttt{overwriteSchema} for deliberate breaking rewrites.
    \item The data files old versions reference; 7 days protects in-flight recovery and time travel.
    \item MERGE is atomic and avoids the delete-then-insert window where readers see missing rows.
\end{enumerate}
\noindent\textbf{Week 3:}
\begin{enumerate}[itemsep=0pt, leftmargin=1.5em]
    \item Bronze is the replayable evidence layer --- every downstream fix starts from raw.
    \item Rows that fail validation, with the rejection reason, so nothing is silently lost.
    \item Partition-scoped \texttt{replaceWhere} and key-based \texttt{MERGE} make reruns converge to the same state.
    \item A star schema gives stable, well-tested dimensions and selective fact scans; wide tables recompute everything.
    \item \texttt{\_ingested\_at}, \texttt{\_source\_file}, and pipeline run IDs.
\end{enumerate}
\noindent\textbf{Week 4:}
\begin{enumerate}[itemsep=0pt, leftmargin=2em]
    \item The driver plans once against the full DAG, enabling predicate pushdown and join reordering.
    \item Wide transformations (groupBy, join) move data across executors --- network and disk cost.
    \item \texttt{repartition} shuffles and can increase or decrease; \texttt{coalesce} only decreases without a full shuffle.
    \item Partition skew --- one task processing far more records than its peers.
    \item When cached data is scanned once, is huge, or forces memory pressure and eviction.
\end{enumerate}
\noindent\textbf{Week 5:}
\begin{enumerate}[itemsep=0pt, leftmargin=2em]
    \item Every file costs listing, footer reads, and a task slot; metadata dominates at small-file scale.
    \item High-cardinality columns used in selective equality/range filters.
    \item The clustering columns --- \texttt{ALTER TABLE ... CLUSTER BY} without rewriting layout decisions into paths.
    \item Rewriting entire Parquet files to remove a handful of deleted rows.
    \item Compaction rewrites files the stream may still be appending; coordinate schedules.
\end{enumerate}
\noindent\textbf{Week 6:}
\begin{enumerate}[itemsep=0pt, leftmargin=2em]
    \item Vectorized execution of scans, joins, and aggregations --- the difference between minutes and seconds on large scans.
    \item When the result is queried repeatedly and source freshness tolerances allow scheduled refresh.
    \item You give up warm-start control and gain instant start, fast autoscaling, and no cluster management.
    \item The warehouse adds clusters up to the max during the burst, then scales back down.
    \item \texttt{system.billing.usage}, joined to \texttt{system.access.audit} or warehouse events.
\end{enumerate}
\noindent\textbf{Week 7:}
\begin{enumerate}[itemsep=0pt, leftmargin=2em]
    \item Genie reasons over semantics; Bronze's raw, uncleaned data produces confident wrong answers.
    \item They steer table/column choice, join paths, and metric definitions for generated SQL.
    \item A curated dataset documents semantics and reuse; direct references couple the dashboard to physical tables.
    \item Certify the Gold view, reconcile it to the source, and lock editing with ownership + lineage.
    \item A warehouse that can run unattended with governed credentials and a refresh schedule.
\end{enumerate}
\noindent\textbf{Week 8:}
\begin{enumerate}[itemsep=0pt, leftmargin=2em]
    \item Metastores bind catalogs, credentials, and audit for a region; typically one per region per cloud.
    \item Managed: Databricks owns file lifecycle (DROP deletes data); external: you own the location.
    \item A view stays live and centrally governed; a copy drifts and multiplies the PII surface.
    \item \texttt{SELECT} on the Gold schema (or specific tables/views) granted to the analyst group.
    \item \texttt{system.access.audit} (with lineage in \texttt{system.access.table\_lineage}/\texttt{column\_lineage}).
\end{enumerate}
\noindent\textbf{Week 9:}
\begin{enumerate}[itemsep=0pt, leftmargin=2em]
    \item Re-plans joins, partitions, and skew handling from runtime statistics instead of compile-time guesses.
    \item When the ``small'' side no longer fits in executor memory, or AQE statistics were wrong.
    \item A Python UDF serializes rows to a Python process per call; pandas UDFs move columnar Arrow batches.
    \item Salting splits the hot key across partitions; the other side must be replicated per salt value.
    \item The stage summary's shuffle read/write columns and the task-duration tail.
\end{enumerate}
\noindent\textbf{Week 10:}
\begin{enumerate}[itemsep=0pt, leftmargin=2em]
    \item Automatic dependency resolution, incrementalization, retries, and quality metrics per dataset.
    \item \texttt{expect\_or\_fail} for PII: bad PII should stop the pipeline, not silently pass or drop.
    \item It records expectations, lineage, and per-run quality --- your evidence base for audits.
    \item When logic or schema changes invalidate incremental state (or for a deliberate rebuild).
    \item Triggered: you pay only for the run; continuous bills idle waiting time.
\end{enumerate}
\noindent\textbf{Week 11:}
\begin{enumerate}[itemsep=0pt, leftmargin=2em]
    \item Notification mode scales to millions of files where directory listing becomes the bottleneck.
    \item The checkpoint records processed files/offsets; restarts resume, never repeat or skip.
    \item Rows that do not match the schema; better than failing because the pipeline survives and you can reprocess.
    \item Type 1 answers ``what is true now''; Type 2 answers ``what was true when''.
    \item It orders CDC events by a sequence/timestamp column so late updates apply in the right order.
\end{enumerate}
\noindent\textbf{Week 12:}
\begin{enumerate}[itemsep=0pt, leftmargin=2em]
    \item Failure isolation: one task's crash or OOM does not strand sibling tasks.
    \item Task values --- set with \texttt{dbutils.jobs.taskValues.set} during a run.
    \item Only the failed and downstream-dependent tasks --- cheaper and faster than a full rerun.
    \item Between Silver validation and Gold publish, as a blocking task.
    \item They route failure paths inside the job graph instead of bolting on external monitors.
\end{enumerate}
\noindent\textbf{Week 13:}
\begin{enumerate}[itemsep=0pt, leftmargin=2em]
    \item It stores offsets and state; without it a restart reprocesses or loses data.
    \item Update mode --- append forbids changing existing aggregate rows.
    \item Batch-like execution of a streaming query: process all available data, then stop.
    \item State schema is checkpointed; incompatible changes fail the restart by design.
    \item Input rate vs.\ processing rate (and the batch duration trend).
\end{enumerate}
\noindent\textbf{Week 14:}
\begin{enumerate}[itemsep=0pt, leftmargin=2em]
    \item State for windows/keys older than the watermark horizon --- bounding memory.
    \item The checkpoint recovers offsets/state; the sink must ignore re-delivered writes (idempotent MERGE/keys).
    \item Overlapping windows --- each event lands in multiple windows.
    \item When state outgrows executor memory and spills must go to local disk/RocksDB.
    \item Deduplication state must remember every seen key forever without a watermark.
\end{enumerate}
\noindent\textbf{Week 15:}
\begin{enumerate}[itemsep=0pt, leftmargin=2em]
    \item Brokers acknowledge writes, consumers track offsets; redelivery is possible --- dedupe at the sink.
    \item Executor overload and memory pressure from oversized micro-batches.
    \item Schemas evolve under contract instead of silently breaking deserialization.
    \item One Kafka partition maps to one Spark task path --- parallelism is partition-bound.
    \item Event Hubs exposes a Kafka-compatible endpoint, so Kafka clients work unchanged.
\end{enumerate}
\noindent\textbf{Week 16:}
\begin{enumerate}[itemsep=0pt, leftmargin=2em]
    \item Streaming DataFrames cannot MERGE directly; \texttt{foreachBatch} runs a batch MERGE per micro-batch with the checkpoint preserving exactly-once.
    \item Delta's optimistic concurrency resolves stream/batch commit conflicts; partition discipline keeps them disjoint.
    \item Sub-second or a few seconds end-to-end; otherwise a 1-minute trigger is far cheaper.
    \item Streams drift from business truth (late data, replays); the batch path is the reconciler of record.
    \item Continuous pays for always-on compute; triggered pays only when processing.
\end{enumerate}
\noindent\textbf{Week 17:}
\begin{enumerate}[itemsep=0pt, leftmargin=2em]
    \item Environment, signatures, input examples, and standard metrics --- the reproducibility scaffolding.
    \item Governance: lineage, access control, and sharing live with the data, not per-workspace.
    \item The alias (e.g.\ \texttt{champion}) --- promotion moves the alias without redeploying jobs.
    \item Logged code version, environment, data snapshot, and seed.
    \item Otherwise you measure data drift, not model quality.
\end{enumerate}
\noindent\textbf{Week 18:}
\begin{enumerate}[itemsep=0pt, leftmargin=2em]
    \item Using information unavailable at prediction time --- e.g.\ ``days\_to\_refund'' in a churn model trained before refund events exist.
    \item The lookup joins each label row only to feature rows with timestamps $\le$ the label's event time.
    \item So every model records exactly which feature tables and versions produced it.
    \item The online store (low-latency lookups); the offline store serves training.
    \item Skew is train-time features computed differently from serve-time; one governed feature pipeline feeds both.
\end{enumerate}
\noindent\textbf{Week 19:}
\begin{enumerate}[itemsep=0pt, leftmargin=2em]
    \item Delta remains the governed source of truth; the index is a derived, rebuildable projection.
    \item Delta-sync --- the index updates automatically as the table changes.
    \item Small chunks improve precision but lose context; large chunks dilute relevance.
    \item Logged metrics make quality regressions visible and comparable across prompt/index changes.
    \item Before embedding/indexing --- never let PII reach the vector store or the model context.
\end{enumerate}
\noindent\textbf{Week 20:}
\begin{enumerate}[itemsep=0pt, leftmargin=2em]
    \item Request/response payloads --- enabling drift and quality monitoring without app changes.
    \item Promotion changes the alias, not the job code --- rollback is re-pointing the alias.
    \item Partition-scoped \texttt{replaceWhere} (or key-based MERGE) so reruns replace rather than duplicate.
    \item Retraining on drift does not guarantee better; promotion needs human or metric gates.
    \item Provisioned concurrency for latency-critical; scale-to-zero for intermittent batch-adjacent calls.
\end{enumerate}
\noindent\textbf{Week 21:}
\begin{enumerate}[itemsep=0pt, leftmargin=2em]
    \item People leave and rotate; service principals outlive individuals and are scoped to the job.
    \item Joiner/mover/leaver group synchronization from the IdP --- no orphaned access.
    \item Key Vault-backed: your cloud vault controls rotation and policy; Databricks-backed is simpler but platform-managed.
    \item Credentials already stolen and used from an allowed network --- it is a perimeter, not authentication.
    \item Data-plane traffic leaves the public internet; exposure shifts to private-endpoint and DNS hygiene.
\end{enumerate}
\noindent\textbf{Week 22:}
\begin{enumerate}[itemsep=0pt, leftmargin=2em]
    \item You enumerate every downstream table/dashboard consuming a column before changing it.
    \item \texttt{system.access.audit}.
    \item Sharing serves data through the protocol's short-lived pre-signed access --- no standing credentials cross the boundary.
    \item Access ends; recipients cannot refresh or re-read the shared data.
    \item Non-Databricks recipients consume via open connectors (pandas, Spark, Power BI) without a Databricks account.
\end{enumerate}
\noindent\textbf{Week 23:}
\begin{enumerate}[itemsep=0pt, leftmargin=2em]
    \item Invalid resource definitions, missing variables, and bad references --- before any cloud change.
    \item Targets isolate workspaces, variables, and deployment state; flags invite accidental prod edits.
    \item Concurrent applies corrupt local state; remote locked state enables team/CI collaboration.
    \item \texttt{system.billing.usage}.
    \item Auto-termination/auto-stop and job clusters (or serverless) instead of always-on all-purpose compute.
\end{enumerate}
\noindent\textbf{Week 24:}
\begin{enumerate}[itemsep=0pt, leftmargin=2em]
    \item Yours to answer --- name the domain and the lab you will rebuild.
    \item Run the clean-clone redeploy drill; a platform that cannot redeploy is not finished.
    \item Cost per month, p95 latency, and rows processed per day.
    \item Yours to answer --- write it in the README's lessons-learned section.
    \item Practice the one-sentence-per-box drill before the review.
\end{enumerate}

% ==========================================
\section{Appendix E: Datasets \& Project Data}
Prefer data already available inside Databricks (\texttt{/databricks-datasets/} in the workspace file store, or the \texttt{samples} catalog) --- it sits next to your compute, so no egress cost or download step.

\begin{itemize}[label=-]
    \item \textbf{NYC Taxi \& Limousine Trip Records} (\texttt{/databricks-datasets/nyctaxi/}, also \url{https://registry.opendata.aws/nyc-tlc-trip-records-pds/}): billions of trip rows in CSV/Parquet. Used in Week 4 Spark tuning, the Week 5 optimization benchmark, and the Milestone 2 analytics capstone.
    \item \textbf{\texttt{samples} catalog} (\texttt{samples.tpch}, \texttt{samples.nyctaxi}): governed sample schemas present in every Unity Catalog workspace --- ideal for SQL warehouse and Genie labs in Weeks 6--7.
    \item \textbf{Olist Brazilian E-Commerce dataset} (\url{https://www.kaggle.com/datasets/olistbr/brazilian-ecommerce}): 100k orders across customers, sellers, products, payments, and reviews --- a natural star schema. Used in the Milestone 1 and Milestone 3 capstones; upload the CSVs to a Unity Catalog volume.
\end{itemize}

\noindent\textbf{Synthetic data:} generate customers, orders, clickstream, and card transactions with Python Faker (\url{https://github.com/joke2k/faker}); the Week 15 Kafka producer script shows the pattern for streaming events. Expected volumes per milestone: Milestone 1 $\sim$1M rows across Bronze--Gold; Milestone 2 $\sim$10M rows optimized Delta; Milestone 3 $\sim$50M raw events per batch; Milestone 4 $\sim$5k events/minute sustained; Milestone 5 $\sim$100k feature rows; Milestone 6 small governed samples ($\sim$100k rows) with injected PII.

% ==========================================
\section{Appendix F: Portfolio \& Presentation Guide}
Each capstone should ship as a public GitHub repository you can walk an interviewer through.

\noindent\textbf{README structure (in this order):}
\begin{enumerate}[itemsep=1pt]
    \item \textbf{Problem:} one paragraph --- who hurts, at what scale, and what ``done'' means.
    \item \textbf{Architecture diagram first:} reviewers decide in 10 seconds; lead with the picture.
    \item \textbf{Setup:} exact commands to reproduce from a clean environment (prereqs, bundle deploy, seed data).
    \item \textbf{Screenshots:} working dashboard, pipeline run, and one alert firing.
    \item \textbf{Cost estimate:} monthly DBU table with the two optimization decisions you made.
    \item \textbf{Lessons learned:} one failure you hit and how you fixed it.
\end{enumerate}

\noindent\textbf{Resume bullet templates (two per milestone --- fill in the brackets):}
\begin{itemize}[label=-, itemsep=1pt]
    \item \textbf{M1:} Built a medallion lakehouse on Databricks (Bronze/Silver/Gold Delta tables) with quarantine accounting reconciling [N] rows/day and time-travel recovery drills.
    \item \textbf{M1:} Tuned PySpark workloads using Spark UI evidence, cutting a [N] TB nightly job from [X] hours to [Y] minutes.
    \item \textbf{M2:} Built a governed analytics layer (liquid-clustered Delta, serverless SQL warehouse, AI/BI + Genie) serving [N] analysts under Unity Catalog PII masking.
    \item \textbf{M2:} Benchmarked Z-ORDER vs.\ liquid clustering, documenting a layout decision saving [X]\% on query latency at [N] rows.
    \item \textbf{M3:} Automated a batch ELT platform with Auto Loader, DLT expectations, and a multi-task Workflow with repair-run recovery, eliminating duplicate loads on re-runs.
    \item \textbf{M3:} Implemented SCD Type 2 CDC with \texttt{AUTO CDC}, preserving full customer history across [N] daily extracts.
    \item \textbf{M4:} Built a real-time fraud-detection stream (Kafka $\rightarrow$ Structured Streaming windows $\rightarrow$ Delta $\rightarrow$ AI/BI) flagging impossible travel within 10 minutes.
    \item \textbf{M4:} Implemented exactly-once streaming with checkpoints and idempotent \texttt{foreachBatch} MERGE, sustaining [N] events/minute with zero duplicate alerts on replay.
    \item \textbf{M5:} Productionized MLOps: point-in-time feature tables $\rightarrow$ MLflow tracking $\rightarrow$ Unity Catalog registry $\rightarrow$ serverless serving with Lakehouse Monitoring drift alarms.
    \item \textbf{M5:} Built a RAG assistant with Mosaic AI Vector Search over [N] documents, evaluating quality with MLflow metrics before launch.
    \item \textbf{M6:} Deployed a governed data platform as code with Databricks Asset Bundles, Terraform, and GitHub Actions dev$\rightarrow$staging$\rightarrow$prod promotion.
    \item \textbf{M6:} Built per-domain chargeback reporting from \texttt{system.billing.usage}, cutting monthly DBU spend [X]\% via auto-termination and serverless policies.
\end{itemize}

\noindent\textbf{Talking about your project in interviews:} open with the problem and its scale (``10M rows nightly, 500 analysts''), then show the architecture diagram and narrate two design decisions with their trade-offs --- interviewers score reasoning, not service names. Have one failure story ready (what broke, how you detected it, what you changed). Close with numbers: DBU cost per month, latency, row counts. If asked ``what would you do differently,'' you pass only if you have a real answer --- write it in the README's lessons-learned section while the project is fresh.

\newpage
% ==========================================
% TRACKING TIMETABLE
% ==========================================
\section{6-Month Progress Tracking Timetable}

\begin{center}
\renewcommand{\arraystretch}{1.5}
\begin{longtable}{@{} c l p{7cm} c @{}}
\toprule
\textbf{Week} & \textbf{Primary Focus} & \textbf{Key Project / Capstone} & \textbf{Status} \\
\midrule
\endfirsthead

\toprule
\textbf{Week} & \textbf{Primary Focus} & \textbf{Key Project / Capstone} & \textbf{Status} \\
\midrule
\endhead

\bottomrule
\endfoot

\bottomrule
\endlastfoot

1 & Databricks Platform \& Workspace & CLI + SDK smoke test, governed compute policy & $\square$ \\
2 & Delta Lake Foundations & Schema enforcement/evolution \& RESTORE drill & $\square$ \\
3 & Medallion Architecture & Bronze/Silver/Gold Olist flow with quarantine & $\square$ \\
4 & Apache Spark on Databricks & \textbf{Cap 1: Retail Medallion Lakehouse} & $\square$ \\
\midrule
5 & Delta Optimization & OPTIMIZE/Z-ORDER vs.\ liquid clustering benchmark & $\square$ \\
6 & Databricks SQL Warehouses & Serverless warehouse + materialized view & $\square$ \\
7 & AI/BI Dashboards \& Genie & Executive dashboard + instructed Genie space & $\square$ \\
8 & Unity Catalog Governance & \textbf{Cap 2: Governed Analytics Lakehouse} & $\square$ \\
\midrule
9 & PySpark Tuning & Skew/shuffle remediation with Spark UI evidence & $\square$ \\
10 & Delta Live Tables & DLT pipeline with expectations \& event log & $\square$ \\
11 & Auto Loader \& CDC & Incremental ingestion + SCD2 customer history & $\square$ \\
12 & Workflows Orchestration & \textbf{Cap 3: Orchestrated Batch ELT Platform} & $\square$ \\
\midrule
13 & Structured Streaming & Checkpointed stream with restart proof & $\square$ \\
14 & State, Watermarks \& Exactly-Once & Watermarked dedup with RocksDB comparison & $\square$ \\
15 & Kafka \& Event Hubs Integration & Kafka$\rightarrow$Delta with replay reconciliation & $\square$ \\
16 & Real-Time Lakehouse & \textbf{Cap 4: Real-Time Fraud Telemetry} & $\square$ \\
\midrule
17 & MLflow \& Tracking & Six autologged runs + champion alias & $\square$ \\
18 & Feature Store & Point-in-time training set with lineage & $\square$ \\
19 & Mosaic AI \& RAG & Vector Search RAG assistant with MLflow eval & $\square$ \\
20 & Serving \& MLOps & \textbf{Cap 5: End-to-End MLOps Pipeline} & $\square$ \\
\midrule
21 & Security, IAM \& Network & Service principal + secret scope hardening & $\square$ \\
22 & Lineage, Audit \& Delta Sharing & Column lineage report + share revocation test & $\square$ \\
23 & Asset Bundles, Terraform \& FinOps & Bundle CI/CD with staging gate + DBU report & $\square$ \\
24 & Final Capstone \& Certification Prep & \textbf{Cap 6: Governed Data Platform} + 2 mock exams & $\square$ \\

\end{longtable}
\end{center}

\vspace{1cm}
\begin{center}
    \textit{"Consistency is the architecture of mastery."} \\
    Best of luck on your 6-month journey to becoming a Master Databricks Data Engineer!
\end{center}

\end{document}
